\documentclass[10pt,a4paper]{article}

% ----------------- Paquetes -------------------%
\usepackage[T1]{fontenc}
\usepackage[utf8]{inputenc}
\usepackage[a4paper,margin=2.5cm]{geometry}
\usepackage{mathptmx}             
\usepackage{changepage} 
\usepackage{setspace}
\usepackage[spanish]{babel}
\usepackage[labelfont=bf,labelsep=space]{caption}
\usepackage{amssymb}
\usepackage{amsmath}
\usepackage{ebgaramond}
\usepackage{placeins}
\usepackage{etoolbox}
\usepackage{setspace}
\usepackage{parskip} 
\usepackage{tcolorbox}
\usepackage{needspace}
\usepackage{xparse}
\usepackage{ocgx2}
\usepackage{hyperref}
\usepackage{hypcap}
\usepackage{soul}\sethlcolor{yellow}% resaltado amarillo: \hl{texto} (Panel TCEE, ⌃H)


%%%%%%%%%%%%%%%%%%%%%%%%%%%%%%%%%%%%%%%%%%%%%%%%%%%%%%%%%%%%%%%%
% --- Fija primero tus valores globales "buenos" ---
\setstretch{1.2}                 % Interlineado global
\setlength{\parskip}{0.7\baselineskip}
\setlength{\parindent}{0pt}
\newlength{\GlobalParskip}
\setlength{\GlobalParskip}{\parskip}

\newcounter{eqnum}[subsection]
\renewcommand{\theeqnum}{\arabic{eqnum}}
\newcounter{alphaitem}
\newcounter{numitem}

%% ------------------- Recuadros----------------------%%
\newcommand{\puntosclave}[1]{%
  \begin{tcolorbox}[
    colframe=black,
    title={Puntos clave},
    before upper={\setlength{\parindent}{0.5cm}\setlength{\parskip}{0.5\baselineskip}}
  ]
  #1
  \end{tcolorbox}
}

\newcommand{\modificaciones}[1]{
  \begin{tcolorbox}[
    colback=white,
    colframe=red,
    title={Modificaciones a realizar},
    before upper={\setlength{\parindent}{0.5cm}\setlength{\parskip}{0.5\baselineskip}}
  ]
  #1
  \end{tcolorbox}
}

%% ------------------- Listas----------------------%%
    % --- Lista alfabética ---
    \newenvironment{la}
      {\begin{list}{\alph{alphaitem})}{%
          \usecounter{alphaitem}%
          \setlength{\leftmargin}{1cm}%
          \setlength{\parskip}{3pt}% 
          \setlength{\itemsep}{0pt}% ← espacio entre ítems
          \setlength{\parsep}{0pt}%  ← párrafos dentro de un ítem
      }}
      {\end{list}}
      
    % --- Lista numérica ---
    \newenvironment{lnum}
      {\begin{list}{\arabic{numitem}.}{%
          \usecounter{numitem}%
          \setlength{\leftmargin}{1cm}%
          \setlength{\parskip}{3pt}% 
          \setlength{\itemsep}{0pt}% ← espacio entre ítems
          \setlength{\parsep}{0pt}%  ← párrafos dentro de un ítem
      }}
      {\end{list}}
      
% ------------------- Imágenes----------------------%
    \graphicspath{{Img/}}% Rutas por defecto 
    % --- Imagen adaptativa ---
    % Registros de dimensión definidos una sola vez
    \newdimen\imgcap
    \newdimen\imgmin
    \newdimen\imgmax
    \newdimen\imgremain
    \newdimen\imgh
    \newdimen\imgneed
    
    \newcommand{\imagenfit}[3]{%
      \addvspace{10pt}% separación antes
      \par\begingroup
        % Parámetros de composición (asignaciones, no \newdimen)
        \imgcap=1\baselineskip            % alto estimado de título+fuente
        \imgmin=0.15\textheight
        \imgmax=0.15\textheight
        \imgremain=\pagegoal
        \ifdim\pagegoal=\maxdimen \imgremain=0pt\fi
        \advance\imgremain by -\pagetotal
        \advance\imgremain by -\imgcap
        % Decisión de altura destino
        \ifdim\imgremain<\imgmin
          \newpage
          \imgh=0.20\textheight
        \else
          \imgh=\imgremain
          \ifdim\imgh>\imgmax \imgh=\imgmax \fi
        \fi
        % Reservar espacio para evitar cortes del bloque completo
        \imgneed=\imgh \advance\imgneed by \imgcap
        \Needspace{\imgneed}
        % Cuerpo
        \noindent\begin{minipage}{\textwidth}
          \centering
          \captionof{figure}{\textemdash\ #2}\par\vspace{0.3em}% título arriba
          \includegraphics[width=\linewidth,height=\imgh,keepaspectratio]{\detokenize{#1}}\par
          \vspace{0.3em}\small\textit{Fuente: }#3% fuente abajo
        \end{minipage}%
      \endgroup
      \par\addvspace{10pt}%
    }

%------------ Bloque de RECUADRO ESPECIAL -------------%
    \newenvironment{specialblock}[1]{%
      \begin{quote} % crea sangría a izquierda y derecha
      \small % texto más pequeño
      \noindent\textbf{#1}\\[-0.3em] % título en negrita
      \rule{\linewidth}{0.4pt}\\[0.6em]}{
      \end{quote}}
      

%-------------------------- Author Font -----------------------%
    \newcommand{\authorfont}[1]{{\fontfamily{EBGaramond-TLF}\selectfont #1}}

%-------------------------- Anexos -----------------------%
    \makeatletter
    \newcounter{anexo}
    \renewcommand{\theanexo}{\Roman{anexo}}
    \newcommand{\anexo}[1]{%
      \clearpage
      \phantomsection
      \refstepcounter{anexo}%
      \noindent\textbf{Anexo \theanexo.- }#1\par\medskip
      \addcontentsline{stc}{section}{Anexo \theanexo.- #1}%
    }
    \makeatother

    %-------------------------- Lista de figuras (personal) -----------------------%
\makeatletter
\newcommand{\MyListOfFigures}{%
  \clearpage
  \phantomsection
  {\centering\bfseries\Large Índice de figuras\par}%
  \medskip
  % Entrada en nuestro índice de contenido (stc)
  \addcontentsline{stc}{section}{Índice de figuras}%
  % Recuperación del listado (sin cabecera propia)
  \begingroup
    \parindent=0pt\parskip=2pt
    \@starttoc{lof}%
  \endgroup
}
\makeatother

%-------------------------- Bibliografía -----------------------%
\makeatletter
% Contador para las referencias
\newcounter{myref}
% Cabecera de la bibliografía, entra en el índice 'stc'
\newcommand{\MyBibliography}{%
  \clearpage
  \phantomsection
  {\centering\bfseries\Large Bibliografía y referencias\par}%
  \medskip
  \addcontentsline{stc}{section}{Bibliografía y referencias}%
  \setcounter{myref}{0}%
}
\newcommand{\MyBibItem}[4]{%
  \refstepcounter{myref}%
  \noindent[\themyref]\ #1.\ \emph{#2}. #3. #4\par
}
\makeatother

%--------- Establecer cuatro niveles de índice --------%
    \setcounter{tocdepth}{4}   
    \setcounter{secnumdepth}{4}% numera hasta \paragraph

%%%%%%%%%%%%%%%%%%%%%%%%%%%%%%%%

\makeatletter

    %--------- Interlineado Global --------%
    \edef\GlobalBaseLineStretch{\baselinestretch}
    
    %--------- Espaciado de seccinones --------%
    \renewcommand\section{\@startsection{section}
      {1}{\z@}
      {2.5ex \@plus 1ex \@minus .2ex} % espacio antes
      {1.5ex \@plus .2ex}             % espacio después
      {\normalfont\Large\bfseries}    % formato del título
    }
    \renewcommand\subsection{\@startsection{subsection}{2}{\z@}
      {3ex \@plus 0.8ex \@minus .2ex}
      {1.5ex \@plus .2ex}
      {\normalfont\large\bfseries}
    }
    \renewcommand\subsubsection{\@startsection{subsubsection}{3}{\z@}
      {3ex \@plus 0.5ex \@minus .2ex}
      {1ex \@plus .2ex}
      {\normalfont\normalsize\bfseries}
    }
    \renewcommand\paragraph{\@startsection{paragraph}{4}{\z@}%
    {-2ex \@plus -0.5ex \@minus -.2ex}% espacio antes
    {0.8ex \@plus .2ex}% espacio después
    {\normalfont\normalsize\itshape}}% ← cursiva en lugar de negrita

    %--------- Ecuaciones --------%   
    % Variables globales para almacenar títulos y notas
    \gdef\leq@current@section{}%
    \gdef\leq@current@subsection{}%
    \gdef\leq@last@printed@section{}%
    \gdef\leq@last@printed@subsection{}%
    
    % Comando para añadir nota (invisible en el cuerpo)
    \newcommand{\eqnote}[1]{%
      \addcontentsline{leq}{leqnote}{#1}%
    }
    
    % Comando para ecuaciones
    \newcommand{\eqblock}[2]{%
      % Verificar si necesitamos imprimir el título de sección
      \ifx\leq@current@section\leq@last@printed@section\else
        \addcontentsline{leq}{leqsection}{\leq@current@section}%
        \global\let\leq@last@printed@section\leq@current@section
        \gdef\leq@last@printed@subsection{}% Reset subsección al cambiar sección
      \fi
      % Verificar si necesitamos imprimir el título de subsección
      \ifx\leq@current@subsection\leq@last@printed@subsection\else
        \ifx\leq@current@subsection\@empty\else
          \addcontentsline{leq}{leqsubsection}{\leq@current@subsection}%
          \global\let\leq@last@printed@subsection\leq@current@subsection
        \fi
      \fi
      % Ecuación
      \refstepcounter{eqnum}%
      \begin{equation}
        #1\tag{E.\theeqnum}
      \end{equation}%
      \addcontentsline{leq}{leqt}{[E.\theeqnum] -- #2}%
      \addcontentsline{leq}{leqm}{#1}%
    }
    
    %%% ----- Índice de ecuaciones ----------%%%
    % Tamaño para []
    \newcommand{\leqIndexFont}{\normalsize}
    \newcommand{\leqTitleFont}{\tiny}
    \newcommand{\leqNoteFont}{\tiny}
    % Cabecera de sección 
    \newcommand\l@leqsection[2]{%
      \addvspace{25pt}% Mayor separación antes de nueva sección
      \noindent\leqIndexFont
      \makebox[\linewidth]{\rule{.38\linewidth}{0.4pt}\ \ \textbf{  \MakeUppercase{#1}}\ \ \rule{.38\linewidth}{0.4pt}}%
      \par\vspace{-10pt}
    }
    % Cabecera de subsección 
    \newcommand\l@leqsubsection[2]{%
      \par\addvspace{15pt}%
      \noindent\leqIndexFont #1
      \par\vspace{-7pt}
    }
    % Título de ecuación 
    \newcommand\l@leqt[2]{%
    \par\addvspace{10pt}%
      \par\noindent\leqTitleFont #1\par
    }
    % Miniatura de ecuación
    \newcommand\l@leqm[2]{%
      \vspace{0.3\baselineskip}%
      \par\scriptsize\noindent\hspace{1.5em}\(#1\)\par%
    }
    % Nota de ecuación 
    \newcommand\l@leqnote[2]{%
      \vspace{0.2\baselineskip}%
      \par\noindent\leqNoteFont\hspace{1.5em}\textbf{Nota.--} #1\par\vspace{0.5\baselineskip}%
    }
    % Generación del índice
    \newcommand{\mylistofequations}{%
      \clearpage
      \phantomsection
      % Título visual de la lista (ajústalo a tu gusto):
      {\centering\bfseries\Large Índice de ecuaciones\par}
      % Entrada en nuestro índice 'stc':
      \addcontentsline{stc}{section}{Índice de ecuaciones}%
      % Llamada a tu lista real (usa el nombre exacto de tu macro):
      \@starttoc{leq}%
    }
    
    %--------- Captura de títulos de sección --------%
    \let\leq@original@section\section
    \let\leq@original@subsection\subsection
    % Flag para saber si estamos en una sección asterisco
    \newif\ifLeqStarSection
    \LeqStarSectionfalse
    
    % Redefinir \section CON CONTROL DE ASTERISCO
    \renewcommand{\section}{%
      \@ifstar{\leq@section@star}{\@dblarg\leq@section@nostar}%
    }
    \def\leq@section@star#1{%
      \global\LeqStarSectiontrue
      \gdef\leq@current@section{#1}%
      \gdef\leq@current@subsection{}%
      \leq@original@section{#1}%
      \global\LeqStarSectionfalse
    }
    \def\leq@section@nostar[#1]#2{%
      \global\LeqStarSectionfalse
      \gdef\leq@current@section{#2}%
      \gdef\leq@current@subsection{}%
      \leq@original@section[#1]{#2}%
    }
    % Redefinir \subsection
    \renewcommand{\subsection}{%
      \@ifstar{\leq@subsection@star}{\@dblarg\leq@subsection@nostar}%
    }
    \def\leq@subsection@star#1{%
      \global\LeqStarSectiontrue
      \gdef\leq@current@subsection{#1}%
      \leq@original@subsection{#1}%
      \global\LeqStarSectionfalse
    }
    \def\leq@subsection@nostar[#1]#2{%
      \global\LeqStarSectionfalse
      \gdef\leq@current@subsection{#2}%
      \leq@original@subsection[#1]{#2}%
    }

    % ----- Restauración de valores Globales de estilo ----------%
    % Flags
    \newif\ifInFootnote
    \InFootnotefalse
    \pretocmd{\@footnotetext}{\InFootnotetrue}{}{}
    \apptocmd{\@footnotetext}{\InFootnotefalse}{}{}
    % Profundidad de listas (soporta anidadas)
    \newcount\MyListDepth \MyListDepth=0
    \AtBeginEnvironment{la}{\advance\MyListDepth by 1}
    \AfterEndEnvironment{la}{\advance\MyListDepth by -1}
    \AtBeginEnvironment{lnum}{\advance\MyListDepth by 1}
    \AfterEndEnvironment{lnum}{\advance\MyListDepth by -1}
    \AtBeginEnvironment{itemize}{\advance\MyListDepth by 1}
    \AfterEndEnvironment{itemize}{\advance\MyListDepth by -1}
    \AtBeginEnvironment{enumerate}{\advance\MyListDepth by 1}
    \AfterEndEnvironment{enumerate}{\advance\MyListDepth by -1}
    % Restauración diferida: hazlo al INICIO del próximo párrafo real
    \newif\if@pendingRestore
    \@pendingRestorefalse
    \newcommand{\RequestRestore}{\global\@pendingRestoretrue}
    \everypar{%
      \if@pendingRestore
        \global\@pendingRestorefalse
        \ifInFootnote\else
          \ifnum\MyListDepth=0
            \setlength{\parskip}{\GlobalParskip}%
            \setstretch{\GlobalBaseLineStretch}%
          \fi
        \fi
      \fi
    }
    % Al final de bloques:
    \AfterEndEnvironment{equation}{\RequestRestore}
    \AfterEndEnvironment{equation}{\RequestRestore}
    \AfterEndEnvironment{la}{\RequestRestore}
    \AfterEndEnvironment{lnum}{\RequestRestore}
    % Al final de macros:
    \apptocmd{\eqblock}{\RequestRestore}{}{}
    \apptocmd{\imagenfit}{\RequestRestore}{}{}
    
\makeatother

% ===================== ToC propio con 4 niveles (único bloque) =====================
\makeatletter

% --- Escritores: añadimos una línea al .stc en cada cabecera numerada ---
% Guardamos las versiones actuales para llamar al comportamiento normal
\let\MyTOC@orig@section\section
\let\MyTOC@orig@subsection\subsection
\let\MyTOC@orig@subsubsection\subsubsection
\let\MyTOC@orig@paragraph\paragraph

% SECTION
\renewcommand{\section}{%
  \@ifstar{\MyTOC@section@star}{\@dblarg\MyTOC@section@nostar}%
}
\def\MyTOC@section@star#1{%
  \MyTOC@orig@section{#1}% sin entrada al .stc
}
\def\MyTOC@section@nostar[#1]#2{%
  \MyTOC@orig@section[{#1}]{#2}% comportamiento normal (escribe al .toc)
  % Escribimos también nuestra línea al .stc (texto con \numberline)
  \addcontentsline{stc}{section}{\protect\numberline{\thesection}#2}%
}

% SUBSECTION
\renewcommand{\subsection}{%
  \@ifstar{\MyTOC@subsection@star}{\@dblarg\MyTOC@subsection@nostar}%
}
\def\MyTOC@subsection@star#1{%
  \MyTOC@orig@subsection{#1}%
}
\def\MyTOC@subsection@nostar[#1]#2{%
  \MyTOC@orig@subsection[{#1}]{#2}%
  \addcontentsline{stc}{subsection}{\protect\numberline{\thesubsection}#2}%
}

% SUBSUBSECTION
\renewcommand{\subsubsection}{%
  \@ifstar{\MyTOC@subsubsection@star}{\@dblarg\MyTOC@subsubsection@nostar}%
}
\def\MyTOC@subsubsection@star#1{%
  \MyTOC@orig@subsubsection{#1}%
}
\def\MyTOC@subsubsection@nostar[#1]#2{%
  \MyTOC@orig@subsubsection[{#1}]{#2}%
  \addcontentsline{stc}{subsubsection}{\protect\numberline{\thesubsubsection}#2}%
}

% PARAGRAPH
\renewcommand{\paragraph}{%
  \@ifstar{\MyTOC@paragraph@star}{\@dblarg\MyTOC@paragraph@nostar}%
}
\def\MyTOC@paragraph@star#1{%
  \MyTOC@orig@paragraph{#1}%
}
\def\MyTOC@paragraph@nostar[#1]#2{%
  \MyTOC@orig@paragraph[{#1}]{#2}%
  % Si no numeras los \paragraph, cambia \theparagraph por texto plano sin \numberline
  \addcontentsline{stc}{paragraph}{\protect\numberline{\theparagraph}#2}%
}

% --- Impresor del índice propio (lee el .stc) ---
\newcommand{\MyTOC}{%
  \par\bigskip
  {\centering\bfseries\Large Índice de contenido\par}%
  \medskip
  \begingroup
    \parindent=0pt\parskip=2pt
    \@starttoc{stc}%
  \endgroup
}

\makeatother
% ===================== fin ToC propio =====================


%%%%%%%%%%%%%%


% --- Datos del tema ---
\title{Sistema Fiscal en España (V): La imposición indirecta. \\
\large Impuesto sobre el Valor Añadido. Impuestos especiales}
\author{Marcelo Ortega}
\date{Septiembre 2026}

\begin{document}
\maketitle
\newpage

% --- Página de resumen y modificaciones ---
\puntosclave{ %% Escribir puntos clave Apuntes CECO
     
    }
\vspace{1cm}

\modificaciones{
    \textbf{NOTA (04/10/2026):} Revisar: https://blog.funcas.es/envejecimiento-de-hogares-nativos-y-recaudacion-de-iva-como-afecta-la-inmigracion/
}

\newpage
\MyTOC
\newpage

\mylistofequations
\clearpage

\MyListOfFigures
\clearpage


\pagenumbering{arabic}
\setcounter{page}{1}


\section{Introducción}



\subsection{Contextualización}


\subsubsection{Relevancia}




\subsubsection{Problemática}



\subsection{Estructura de la exposición}







\clearpage
\section{Impuesto sobre el Valor Añadido: Marco jurídico-conceptual}
\puntosclave{
    
}

España no llega al IVA en 1986 desde el vacío: llega arrastrando seis siglos y medio de intentos, fracasos y reformas fallidas de gravar el consumo y las transacciones, una historia propia que merece el mismo detenimiento que ya dedicamos a la historia general. Este primer bloque recorre cuatro etapas de esa historia específicamente española —los antecedentes premodernos, la larga agonía del impuesto de consumos en el siglo XIX y principios del XX, el tributo en cascada del franquismo, y el nacimiento del IVA español propiamente dicho—, antes de pasar al marco jurídico vigente (punto 2), la gestión del impuesto (punto 3) y el esquema básico de liquidación (punto 4). No me preocupa duplicar contenido ya trabajado en el Tema 13: aquí el foco es exclusivamente español, y cuando la historia española converja con la europea —la adhesión de 1986—, lo señalaré como punto de conexión, no como repetición.


\subsection{Desarrollo histórico: ¿Cómo hemos llegado hasta aquí?}
La figura fiscal más antigua e importante en España (Castilla del Antiguo Régimen) y, por tanto, el antecedente español más directo de un impuesto general sobre las transacciones, es la alcabala, cuyo nombre procede del árabe hispánico \textit{al-qabāla} (el contrato). 

Nace en las Cortes de Burgos de 1342, bajo Alfonso XI, como un servicio extraordinario votado para financiar la campaña militar de Algeciras dentro de la Reconquista, técnicamente heredero de la vieja \textit{vicesima rerum venalium} romana (gravamen del veinteavo sobre cualquier compraventa o permuta, [\authorfont{Ver Tema 4.B.13}]). Nacida como tributo concejil de carácter local, se generalizó bajo Alfonso XI y se prorrogó sucesivamente en Cortes hasta consolidarse, ya de forma permanente, como impuesto ordinario del reino. En 1349 el tipo se duplicó al $10\%$, porcentaje que se mantuvo vigente durante buena parte de los tres siglos siguientes. 

Como podemos intuir, el problema económico de la alcabala es el mismo que cualquier figura tributaria indirecta clásica: al gravar cada transacción sucesiva de un mismo bien era, de facto, un impuesto en cascada ([\authorfont{Ver Tema 4.B.13}]). 

Sobre la alcabala se fueron superponiendo, a lo largo de los siglos XVI y XVII, nuevas figuras que agravaban el mismo problema. Los cientos (\textit{cuatro unos por cien}) eran una ampliación directa de la alcabala, un recargo adicional del $4\%$ sobre las mismas operaciones ya gravadas. Los millones, subsidios de Cortes originalmente extraordinarios aprobados bajo Felipe II a finales del siglo XVI y que terminaron convirtiéndose en permanentes, cargaban específicamente sobre la venta al por menor de artículos de consumo básico: vino, vinagre, aceite, carne, jabón, velas de sebo (precedente directo del debate sobre la diferenciación de tipos). Junto a estas figuras convivían las tercias reales (participación de la Corona en el diezmo eclesiástico sobre granos, vinos y menudos) y un conjunto disperso de rentas menores (martiniega, marzazga, fumazga, entre otras). Al conjunto de estas figuras se le dio, desde comienzos del siglo XVIII, el nombre colectivo de rentas provinciales, por gravar exclusivamente a las provincias de Castilla.

El sistema de recaudación de estas rentas provinciales merece mención propia porque anticipa un problema que reaparecerá en la historia española del siglo XIX. Se gestionaban mayoritariamente mediante arrendamiento (\textit{encabezamientos generales del reino}, perpetuados desde 1536), es decir, el Estado cedía a particulares el derecho a recaudar a cambio de una cantidad fija, quedándose el arrendatario con la diferencia entre lo recaudado y lo entregado a la Corona. El resultado, documentado por la historiografía económica, fue un sistema que gravaba con mucha más intensidad el consumo cotidiano de bienes básicos que la riqueza inmobiliaria, generando beneficios muy sustanciales para los arrendatarios y un descontento social profundo y sostenido en el tiempo.

El primer intento serio de sustituir por completo este entramado por una fiscalidad más racional llega con el proyecto de única contribución del Marqués de la Ensenada, bajo Fernando VI, a mediados del siglo XVIII. Sustentado en el ambicioso Catastro de Ensenada (1749-1756), un censo exhaustivo de la riqueza de las $22$ provincias de la Corona de Castilla, tenía el objetivo de sustituir todas las rentas provinciales (alcabala, cientos,, etc.) por un único impuesto directo sobre la riqueza real de cada contribuyente. El proyecto, pese a su ambición y al ingente trabajo administrativo que generó, nunca llegó a implementarse en su forma plena, y las rentas provinciales sobrevivieron, con reformas parciales sucesivas, hasta bien entrado el siglo XIX.

\subsubsection{El siglo XIX y principios del XX: la larga agonía del impuesto de consumos}
El heredero directo de la alcabala en la España del siglo XIX es el impuesto de consumos, cuya trayectoria es, probablemente, la historia de mayor inestabilidad política de toda la fiscalidad española: abolido en 1813 por las Cortes de Cádiz, en pleno auge liberal; restablecido en 1816 tras la vuelta al absolutismo de Fernando VII; suprimido de nuevo durante el Trienio Liberal (1820-1823); y reinstaurado en 1824. Esta secuencia revela ya un patrón que se repetirá durante todo el siglo. El impuesto de consumos se convierte en un marcador ideológico entre el liberalismo progresista, que lo suprime cada vez que llega al poder, y el liberalismo moderado, que lo reinstaura.

La forma que el impuesto conservaría durante el resto del siglo llega con la reforma de Alejandro Mon de 1845. Este crea la contribución de consumos propiamente dicha, gravando bienes específicos y conservando los antiguos \textit{derechos de puertas} ([\authorfont{Ver Tema 4.B.13}]). No obstante, el impuesto siguió fluctuando durante el resto del siglo. Suprimido en 1868 tras la Revolución que derrocó a Isabel II, lo que dejó a los ayuntamientos españoles sin su principal fuente de ingresos y contribuyó, en varias provincias, al éxito inicial del levantamiento cantonalista de 1873. Restablecido después, y suprimido de nuevo, también por razones de legitimidad política, durante la Primera República. 

En definitiva, durante el periodo de la Restauración, el impuesto de consumos se convirtió en uno de los pilares documentados del caciquismo local. Su gestión municipal, con frecuencia arrendada a particulares, generaba una red de favores, exenciones discrecionales y arbitrariedad recaudatoria que la historiografía sobre provincias ha identificado como una de las bases materiales del control político de los caciques sobre la vida municipal española de la época.

El impuesto fue finalmente abolido por ley de 12 de junio de 1911, bajo el reinado de Alfonso XIII, y sustituido por arbitrios municipales alternativos. No obstante, su desaparición no fue del todo definitiva ni inmediata, ya que la propia ley preveía un calendario escalonado según la vigencia de los contratos de arrendamiento existentes en cada municipio, y determinadas formas de imposición indirecta local sobre el consumo sobrevivieron en la práctica durante años. 

Durante la dictadura de Primo de Rivera (1923-1930), FLORES DE LEMUS proyectó, dentro de una reforma más amplia de la hacienda local, la supresión definitiva de cualquier resto del viejo impuesto de consumos, cerrando así, más de un siglo después de la reforma de MON, un capítulo que había atravesado prácticamente todos los regímenes políticos españoles del siglo XIX y principios del XX, sin que ninguno de ellos consiguiera resolverlo de forma estable.

\subsubsection{El franquismo y la Reforma Tributaria de 1964: el Impuesto General sobre el Tráfico de las Empresas}
El sucesor moderno de todo este linaje de impuestos sobre transacciones llega con la Reforma Tributaria de 1964, impulsada por el ministro de Hacienda NAVARRO RUBIO (la misma reforma que, en su vertiente de imposición directa, ya trabajamos [\authorfont{Ver Tema 4.B.12}]). 

La reforma articuló el sistema en torno a tres figuras indirectas centrales: el Impuesto sobre Transmisiones Patrimoniales y Actos Jurídicos Documentados, el Impuesto General sobre el Tráfico de las Empresas (IGTE o ITE) y el Impuesto de Lujo; completadas por los Impuestos Especiales sobre el Consumo y la Renta de Aduanas.\footnote{%
    La propuesta de gravar el volumen de ventas mediante un impuesto de esta naturaleza había sido reclamada, ya antes de la reforma, por economistas como FUENTES QUINTANA.
}

El ITE era, técnicamente, el misma tipo de impuesto que indentificaríamos con el nombre de tributo en cascada. Gravaba el importe total de la contraprestación en cada una de las transacciones de bienes o servicios realizadas en cualquier fase del proceso de producción o distribución, sin mecanismo de deducción alguno del impuesto ya soportado en la fase anterior.\footnote{
No lo tratamos en profundidad aquí ([\authorfont{Ver Tema 4.B.14}]), pero conviene conocer las críticas que la propia exposición de motivos de la ley española del IVA de 1985 dirigiría contra él. Fundamentalmente tres:
\begin{lnum}
    \item \textbf{Distorsiones organizativas: producción}. La acusada falta de neutralidad interior del ITE constituía un importante factor de interferencia en la correcta asignación de los recursos económicos que impulsa a la utilización de ciclos productivos cortos y a la concentración empresarial. Es decir, un incentivo a la integración vertical artificial;
    \item \textbf{Distorsiones externas: exportación y ajustes}. También criticaba que, desde la perspectiva del comercio exterior, la falta de neutralidad y transparencia de los impuestos en cascada había determinado la desaparición de esta modalidad de impuesto en el área territorial de la Comunidad Económica Europea, precisamente porque la imposibilidad de conocer con precisión la carga tributaria acumulada soportada por un producto exportado generaba una deficiente compensación a dichos exportadores o el establecimiento de eventuales y no deseadas primas a la exportación. Es decir, un problema de ajuste fiscal en frontera; y,
    \item \textbf{Distorsiones productivas: ineficiencias}. Por último, que, al carecer de mecanismo de deducción, el impuesto incidía de manera más intensa sobre los bienes de equipo adquiridos por la empresa, los cuales resultan doblemente gravados en su adquisición y en la enajenación de los productos finales del proceso, la misma penalización de la inversión productiva. Es decir, una violación directa del principio de eficiencia productiva que DIAMOND y MIRRLEES formalizarían apenas unos años después. 
\end{lnum}
}

Además, un dato señalado por la doctrina revela un problema práctico de control. Dado que el consumidor final percibía, formalmente, un único ITE en el precio de lo que compraba, mientras que el producto podía arrastrar tantas cargas de ITE ocultas como fases intermedias hubiera atravesado, era relativamente habitual que las transacciones entre mayorista y minorista no llegaran siquiera a declarar correctamente el impuesto correspondiente a la Administración de la época. Por tanto, persistía un problema de opacidad y de verificabilidad que el propio mecanismo de crédito de factura del IVA, con su cadena de autocontrol documental entre empresas, resolvería después de raíz.

\subsubsection{La adhesión a la CEE y el nacimiento del IVA español: la Ley 30/1985}
España, que había solicitado formalmente su adhesión a la Comunidad Económica Europea en 1977, necesitaba adoptar el sistema común de IVA como condición de esa adhesión. 

La propia exposición de motivos de la Ley 30/1985, de 2 de agosto, del IVA, lo formula con precisión: \textit{la reforma de la imposición indirecta española viene exigida no sólo por imperativos de nuestra futura adhesión a la Comunidad Económica Europea, sino también por razones objetivas de indudable relevancia}. Reconociendo explícitamente que la sustitución del ITE no era solo un trámite de adhesión, sino una reforma que España necesitaba por sus propios méritos técnicos, exactamente los problemas que la figura anterior generaba. La misma presión externa y necesidad interna genuina, que ya identificamos en la historia general del IVA europeo con el caso de LAURÉ y la crisis del impuesto en cascada francés de los años cincuenta.

La Ley 30/1985 entró en vigor el 1 de enero de 1986, coincidiendo exactamente con la fecha de adhesión efectiva de España a la Comunidad Económica Europea. Pero, en un detalle que conviene no pasar por alto, el ITE no desapareció por completo.

La propia disposición final de la ley del IVA lo derogó solo parcialmente, y el impuesto sobrevivió en los territorios excluidos de la armonización comunitaria del IVA (Canarias (IGIC), y Ceuta y Melilla).\footnote{%
    Además, hay un dato de sorprendente actualidad que conviene retener como cierre de este epígrafe, porque conecta esta historia de 1986 con la política fiscal española contemporánea: como consecuencia de la desaparición del ITE en Canarias, el Estado venía abonando anualmente al archipiélago unas compensaciones económicas, en reconocimiento del impacto fiscal de aquella sustitución normativa de hace cuatro décadas — compensaciones que el Consejo de Ministros no acordó suspender hasta diciembre de 2019, es decir, treinta y tres años después de que el IVA sustituyera formalmente al viejo impuesto en cascada franquista.
}

\subsubsection{De 1985 a hoy: la Ley 37/1992 y la evolución de los tipos}
La Ley 30/1985 fue sustituida por la actual Ley 37/1992, precisamente para adaptar el IVA español al régimen del Mercado Único europeo, vigente desde el 1 de enero de 1993 — el mismo régimen transitorio de tributación en destino para las operaciones intracomunitarias, nacido como parche provisional ante la imposibilidad de implantar el sistema de origen que Europa se había propuesto originalmente, y que sigue, más de tres décadas después, sin sustituirse por el sistema definitivo que en su día se prometió. 

Es precisamente esta Ley 37/1992, junto con su reglamento aprobado por Real Decreto 1624/1992, la que rige hoy el IVA español y que desarrollaremos con detalle.

La evolución del tipo general de IVA desde su implantación merece un repaso propio, porque traza la historia fiscal reciente de España condensada en una sola cifra. El tipo general nació, en 1986, en el 12\%; se elevó al 13\% con la entrada en vigor de la Ley 37/1992; alcanzó el 15\% en 1995; y permaneció estable, durante más de una década, en el 16\%.  — el punto de partida exacto que el propio documento ICEX-CECO menciona [ICEX-CECO] al señalar que "el tipo general IVA en un periodo corto de tiempo y a través de dos subidas pasó del 16% al 21%"—.

Esas dos subidas: la elevación al 18\% en julio de 2010, bajo el Gobierno de Zapatero, en el contexto de la crisis financiera internacional y el primer gran ajuste fiscal de la crisis española; y la elevación al 21\%, el tipo actualmente vigente, en septiembre de 2012, bajo el Gobierno de Rajoy, en máxima presión de los mercados sobre la deuda soberana española. 

Esta misma tendencia alcanzó a los tipos reducidos y a buena parte de los tipos de los Impuestos Especiales, respondiendo a las recomendaciones de organismos internacionales favorables a desplazar hacia la imposición indirecta la carga que recae sobre la imposición directa y las cotizaciones sociales: ensanchar la base indirecta, aliviar la directa. 

El resto de este bloque parte ya del marco jurídico vigente, sin necesidad de volver a justificar por qué España llegó hasta él.


\subsection{Marco jurídico actual}
El IVA está regulado actualmente por la Ley 37/1992, de 28 de diciembre, cuya génesis histórica ya hemos trazado. La desarrolla el Real Decreto 1624/1992, de 29 de diciembre, que aprueba su Reglamento. 

En el ámbito comunitario, la norma de referencia es la Directiva 2006/112/CE del Consejo, de 28 de noviembre de 2006, que codificó y sustituyó a la Sexta Directiva de 1977 sin alterar su contenido esencial. Se aplica también teniendo en cuenta los regímenes de Concierto y Convenio Económico vigentes en el País Vasco y Navarra, así como los tratados y convenios internacionales incorporados al ordenamiento español.

El territorio de aplicación del impuesto (TAI) es el territorio español, incluidas las islas adyacentes, el mar territorial hasta las 12 millas náuticas y el espacio aéreo correspondiente (con la exclusión de Canarias, IGIC, y de Ceuta y Melilla, donde rige el Impuesto sobre la Producción, los Servicios y la Importación, IPSI). 

\subsubsection{Naturaleza y características del impuesto}
El IVA es un tributo: (i) \textbf{multifásico}, lo que permite un mayor control administrativo, precisamente por el mecanismo de autocontrol documental entre empresas; (ii) \textbf{no acumulativo}, con menores efectos sobre el nivel general de precios que un impuesto en cascada como el ITE; (iii) \textbf{neutral}, tanto a nivel interior como exterior; (iv) \textbf{general}, porque grava la totalidad de los procesos de producción y distribución; (v) \textbf{indirecto}, al gravar una manifestación indirecta de capacidad contributiva; (v) \textbf{objetivo}, porque no atiende a las circunstancias personales del sujeto pasivo; (vi) \textbf{instantáneo}, devengándose operación por operación; y, (vii) \textbf{parcialmente cedido} a las CCAA ($50\%$), de acuerdo con el índice de consumo territorial certificado por el INE.

\subsubsection{Hecho imponible: Territorio y presupuestos gravados}
La Ley grava tres hechos imponibles distintos: (i) las entregas de bienes y prestaciones de servicios efectuadas por empresarios o profesionales; (ii) las adquisiciones intracomunitarias de bienes; y, (iii) las importaciones de bienes. 

\paragraph{Convencional: Entrega de bienes y servicios}
Para definir el hecho imponible más \textit{general}, en el sentido de que resulta el ejemplo clásico al que recurrir cuando se intenta definir el IVA, necesitamos resolver antes dos cuestiones previas:
\begin{la}
    \item \textbf{¿Qué es un empresario?}. La Ley define como empresarios o profesionales, entre otros: (i) quienes realicen actividades empresariales o profesionales (salvo a título exclusivamente gratuito); y, (ii) las sociedades mercantiles, salvo prueba en contrario.\footnote{%
        La lista es notablemente más extensa, incluyendo, entre otros: quienes exploten un bien corporal o incorporal con el fin de obtener ingresos continuados en el tiempo, en particular los arrendadores; quienes urbanicen terrenos o promuevan/construyan/rehabiliten edificaciones destinadas a su venta, adjudicación o cesión, aunque sea ocasionalmente; y los particulares que entreguen ocasionalmente un medio de transporte nuevo exento por destinarse a otro Estado miembro. Además, se presume el ejercicio de actividad empresarial o profesional cuando se anuncie un establecimiento por cualquier medio publicitario, o cuando las operaciones exijan contribuir por el IAE.
    } Aunque la lista no sea exhaustiva, la intuición prevalece: ¿por qué esta definición es, a priori, tan compleja o poco general, como pasa con otros impuestos? En realidad, lo que a primera vista parece un simple ejercicio de delimitación, tiene una función estructural en la arquitectura del impuesto.\footnote{%
        Como sabemos, según el Teorema DIAMOND y MIRRLEES ([\authorfont{Ver Tema 4.B.13}]) la neutralidad del sistema exige que solo la transacción con el consumidor final soporte carga efectiva, permaneciendo neutras todas las transacciones intermedias entre empresas. Por tanto, es imperativo que la ley sepa identificar con precisión quién es empresa a estos efectos y quién es consumidor final.
    } Si el concepto de empresario o profesional fuera demasiado estrecho, ciertas transacciones que en realidad son intermedias quedarían fuera del mecanismo de crédito de factura y resurgiría el problema de acumulación piramidal, en miniatura, precisamente en ese punto de la cadena. Por tanto, la amplitud deliberada es la condición técnica para que el mecanismo de neutralidad funcione sin fisuras;
    \item \textbf{¿Qué es una entrega de bienes y prestaciones de servicios?}. Se entiende por entrega de bienes la \textbf{transmisión del poder de disposición sobre bienes corporales} (incluidos el gas, el calor, el frío y la energía eléctrica).\footnote{%
        La norma considera también entregas de bienes, aunque no impliquen técnicamente transmisión del poder de disposición, ciertos servicios concretos. Por ejemplo, las ejecuciones de obra de construcción o rehabilitación cuando quien ejecuta la obra aporta materiales cuyo coste supere el $40\%$ de la base imponible, la transmisión de bienes por resolución administrativa o judicial, incluida la expropiación forzosa, las transmisiones entre comitente y comisionista que actúa en nombre propio, y el suministro de software normalizado en soporte material. 
    
    Esta regla del $40\%$ responde, de nuevo, a una lógica económica precisa. Distingue, dentro de un mismo contrato de ejecución de obra, entre aquellos en los que predomina el valor del material aportado (que se asemejan económicamente a una compraventa de bien) y aquellos en los que predomina el valor del trabajo prestado (que se asemejan a una prestación de servicios pura), una distinción con consecuencias reales sobre las reglas de localización territorial y de devengo.
    } Por otro lado, las \textbf{prestaciones de servicios son, por definición residual}, toda operación sujeta que no sea entrega, adquisición intracomunitaria o importación.\footnote{%
        Entre otros: el ejercicio independiente de profesión, arte u oficio; arrendamientos de bienes, industria o negocio; cesiones de uso o disfrute; cesiones de propiedad intelectual e industrial; ejecuciones de obra que no sean entregas; transportes; hostelería y restauración; seguros; hospitalización; préstamos y créditos; y una larga lista adicional ya recogida en el documento de referencia.
    } Además, existe un caso particular que la Ley contempla: los \textbf{autoconsumos}, que la ley asimila a entregas de bienes o a prestaciones de servicios, pese a no mediar contraprestación alguna. Entran dentro de esta definición otras operaciones, como son la transferencia de bienes del patrimonio empresarial al personal, el cambio de afectación entre sectores diferenciados de actividad, o la aplicación de bienes o servicios a fines particulares.\footnote{%
        Conviene no limitarse a listar sino explicar el porqué teórico. Si un empresario pudiera retirar libremente bienes de su patrimonio empresarial para consumo propio sin que ello generara ningún efecto fiscal, se abriría una vía de escape sistemática frente al principio de que toda la cadena debe converger, tarde o temprano, en un acto de consumo final gravado. Pensemoslo, el empresario podría autoabastecerse sin soportar nunca el IVA que un consumidor ordinario sí soporta al comprar el mismo bien en el mercado. Por tanto, la figura del autoconsumo es el parche técnico que cierra esa vía de escape, garantizando que ningún bien o servicio salga del circuito empresarial hacia el consumo final (propio o ajeno) sin generar el gravamen que la neutralidad del sistema exige.
    }
\end{la}

Cuando concurran estas dos circunstancias, el hecho imponible será gravado de acuerdo con las normas que luego expondremos. Ahora bien, existen algunas situaciones, tasadas por Ley, en las que, aún concurriendo dichas condiciones, la operación no está sujeta a gravamen, esta exenta. De entre estas, las más relevantes son las entregas gratuitas de muestras y las prestaciones de demostración con fines de promoción, los servicios prestados por trabajadores por cuenta ajena o las entregas de dinero a título de contraprestación o pago.\footnote{%
    También, la transmisión de un conjunto de elementos que constituya una unidad económica autónoma dentro del patrimonio empresarial, o las entregas y prestaciones realizadas directamente por entes públicos sin contraprestación (con las excepciones ya señaladas cuando actúan mediante empresa mercantil, o en sectores como telecomunicaciones, suministro de agua/gas/electricidad o transporte).
}

\paragraph{Operaciones internacionales: intracomunitarias y extracomunitarias}
Como ya adelantamos, cuando la Comunidad Económica pretendió crear un mercado interior sin fronteras fiscales aspiraba a gravar en origen. No obstante, la imposibilidad práctica de esa aspiración obligó a mantener el régimen transitorio de tributación en destino, todavía vigente hoy. 

La legislación española del IVA implementa este régimen mediante un hecho imponible propio: la adquisición intracomunitaria de bienes (AIB). El mecanismo, aplicado ahora a la letra concreta de la LIVA, reproduce exactamente el esquema que ya hemos expuesto. Cuando un empresario español compra mercancía a un proveedor de otro Estado miembro, esa entrega queda exenta en el país de origen (el vendedor no repercute IVA de su propio país), y el comprador español debe autoliquidar el IVA español correspondiente a esa adquisición, declarándolo simultáneamente como devengado y como deducible en la misma autoliquidación si tiene derecho a deducción plena (mecanismo de autorrepercusión). 

La ley distingue dos supuestos de AIB: las realizadas a título oneroso por empresarios/profesionales o personas jurídicas que no actúen como tales, cuando el transmitente sea empresario/profesional; y las de medios de transporte nuevos, cualquiera que sea el destinatario. Quedan no sujetas las AIB realizadas por sujetos acogidos al régimen especial de agricultura/ganadería/pesca, por quienes realicen exclusivamente operaciones sin derecho a deducción, y por personas jurídicas que no actúen como empresarios.\footnote{%
    Salvo que el importe total de sus adquisiciones a otros Estados miembros supere, en el año natural precedente, los 10.000 euros, umbral por debajo del cual se presume que su volumen de comercio transfronterizo es tan reducido que no justifica activar el mecanismo completo de autoliquidación,  umbral de proporcionalidad administrativa del mismo tipo, aunque de naturaleza distinta, que el que trabajamos a propósito de los umbrales de registro para pequeños negocios ([\authorfont{Ver Tema 4.B.13}]).
}

Las importaciones de bienes procedentes de terceros países siguen un mecanismo mucho más simple, precisamente porque aquí sí existe una frontera aduanera real y verificable. Cualquier entrada de un bien procedente de un territorio tercero está sujeta al impuesto, con la única complejidad añadida de determinar correctamente la base imponible, sin necesidad de recurrir al mecanismo de exención-más-autoliquidación que sí hace falta para las AIB.\footnote{%
    Se calcularía como el valor en aduana más derechos, exacciones y gastos accesorios hasta el primer lugar de destino en la Comunidad.
} En esencia, la frontera aduanera exterior sirve como punto de control físico del bien, por tanto basta por sí solo para garantizar el gravamen.

\subsubsection{Base imponible: Materialización del impuesto}
Conocer el hecho imponible, el conjunto de operaciones sujetas a gravamen, no es suficiente para la implementación efectiva del impuesto. Sería únicamente determinar qué es lo que queremos gravar, pero dejando de lado cuestiones esenciales como el cuándo y el cómo gravarlo. 

\paragraph{Territorio de Aplicación del Impuesto (TAI): ¿dónde se grava la operación?}
En primer lugar, debemos determinar cuál es el alcance de la potestad tributaria del Estado. Naturalmente, esta se restringe a los límites territoriales (con notables excepciones), por tanto la operación estará siempre, y únicamente, sujeja a gravamen cuando la localización de la operación esté dentro del territorio de aplicación del impuesto (TAI).

La complejidad, obviamente, reside en el comercio internacional de servicios. ¿Por qué? Porque en el comercio de mercancías pueden suceder dos únicas situaciones: o bien entra por la frontera exterior, y por tanto tributa en aduanas, o bien es objeto de tráfico interior, por lo que está (objetivamente) dentro del TAI. Para la entrega de bienes, en consecuencia, la regla general será considerar que la operación ha sido realizada en el TAI cuando se pongan a disposición del adquirente en dicho territorio, siempre y cuando no sean objeto de expedición o transporte.\footnote{%
    Se entienden también realizadas en el TAI: las entregas de bienes muebles corporales objeto de expedición o transporte, cuando esta se inicie en el propio territorio (con independencia de dónde finalice, dentro del ámbito comunitario); las entregas de bienes que deban instalarse o montarse, siempre que impliquen su inmovilización; las entregas de bienes inmuebles que radiquen en el territorio; y las entregas realizadas a pasajeros a bordo de un buque, avión o tren, por la parte del trayecto que discurra dentro de la Unión, cuando el origen esté en el TAI.
} 

Por otro lado, se considerará que las operaciones de prestación de servicios han sido realizadas en el TAI siempre que se cumpla, alternativamente: (i) o bien que el servicio se localiza donde radique el destinatario (cliente) o su actividad económica, cuando sea prestado por un empresario o profesional, con independencia de su residencia o desde el lugar del cual preste el servicio (regla B2B); o, (ii) bien que el servicio se localiza donde radique la residencia del prestador, o su actividad económica, cuando el destinatario no actúe como empresario o profesional (la regla B2C).\footnote{%
    Existen, además, reglas especiales para supuestos concretos: servicios relacionados con inmuebles (localizados donde radique el inmueble), servicios electrónicos prestados a consumidores finales (localizados donde resida el destinatario) y ejecuciones de obra sobre bienes muebles corporales (localizadas donde se presten materialmente).
}

\paragraph{Devengo: ¿cuándo se grava la operación?}
Una vez sabemos que la operación se ha realizado dentro del TAI, debemos determinar si se ha materializado o no a efectos de la obligación. Es decir, si concurren la circunstancia temporal para que sea objeto de tributación. 

El devengo, el momento exacto en que nace la obligación tributaria, sigue también una regla dual: (i) en las entregas de bienes, se devenga cuando tiene lugar la puesta a disposición del adquirente; y, (ii) en las prestaciones de servicios, cuando se presten, ejecuten o efectúen efectivamente las operaciones gravadas.

\paragraph{Valor: ¿cuánto está sujeto a tributación?}
Por último, debemos determinar cuál es el monto total de la operación. Es decir, cuál será la Base Imponible o el valor total del negocio sujeto a la obligación tributaria. 

En este sentido, la base imponible está constituida por el importe total de la contraprestación, tanto si se trata de entregas de bienes como si es una prestación de servicios.\footnote{%
    Se incluyen en ese importe: los gastos de transporte, seguros, portes y comisiones vinculados a la operación principal; las subvenciones directamente vinculadas al precio; y los envases y embalajes, incluso los susceptibles de devolución. No se incluyen: las indemnizaciones percibidas; los descuentos y bonificaciones concedidos previa o simultáneamente a la realización de la operación; y los suplidos, cantidades pagadas en nombre y por cuenta del cliente, que no forman parte de la contraprestación propia del sujeto pasivo.
}

\subsubsection{Sujeto Pasivo: Responsable de la obligación}
La regla general es clara: son sujetos pasivos los empresarios o profesionales que realizan las entregas de bienes y prestaciones de servicios sujetas al impuesto. 

\paragraph{Inversión del sujeto pasivo: ¿solución temporal al fraude del operador desconocido?}
Lo que corresponde aquí es desarrollar cómo responde concretamente el ordenamiento español, más allá de la regla general de inversión del sujeto pasivo frente a no residentes. El art. 84 LIVA no se limita a la regla general de inversión del sujeto pasivo frente a operaciones con no residentes, sino que contiene un catálogo de supuestos de inversión doméstica, introducidos sucesivamente como respuesta a patrones de fraude identificados en sectores concretos.\footnote{%
    Entre ellos: las entregas de derechos de emisión y sus certificados; determinadas ejecuciones de obra en el sector de la construcción, con y sin aportación de materiales, en el marco de contratos directos entre promotor y contratista principal cuando tengan por objeto la urbanización de terrenos o la construcción/rehabilitación de edificaciones; las entregas de inmuebles en determinados supuestos, incluidas las producidas en procesos concursales; las entregas de gas y electricidad a determinados empresarios revendedores; y, desde el 1 de abril de 2015, las entregas de plata, platino y paladio en bruto o semilabrado, y muy señaladamente las de teléfonos móviles, consolas de videojuegos, ordenadores portátiles y tabletas digitales.
}

El diseño de estas extensiones es un ejemplo perfecto de cómo el legislador debe calibrar el alcance de la inversión del sujeto pasivo para no generalizarla más allá de lo estrictamente necesario.\footnote{%
    Si extendieramos este mecanismo a toda la cadena, desplazaríamos el riesgo de fraude hacia el último eslabón en lugar de eliminarlo
} En estos productos concretos, la inversión se activa bajo dos condiciones alternativas: (a) cuando el destinatario es un empresario o profesional que se dedica habitualmente a la reventa de estos bienes, debiendo comunicar previamente esa condición de revendedor a la Administración mediante declaración censal antes de iniciar esa actividad; o, (b) cuando el destinatario no es revendedor habitual, pero el importe total de las entregas de estos bienes documentadas en una misma factura supera los $10.000$ euros, excluido el IVA. 

Por debajo de ese umbral, o fuera de esa condición de revendedor, rige la regla general: el vendedor repercute normalmente. Las operaciones sujetas a esta inversión deben además documentarse en una serie de facturación específica, separada del resto de operaciones del sujeto pasivo. Un requisito formal que permite a la Administración aislar y monitorizar con precisión el volumen de estas transacciones de alto riesgo, sin tener que rastrearlas entre el resto del tráfico ordinario de la empresa.

\paragraph{Suministro Inmediato de Información: ¿la solución definitiva al fraude de operador desconocido?}
Junto a esta batería de medidas dirigidas a transacciones concretas, España desarrolló de forma pionera, un instrumento de control complementario que no actúa reasignando la obligación de pago sino acelerando la verificación: el Suministro Inmediato de Información (SII). 

El SII obliga a remitir electrónicamente a la Agencia Tributaria, en un plazo de apenas unos días desde su expedición o recepción, el contenido íntegro de los libros registro de facturas emitidas y recibidas, de bienes de inversión y de determinadas operaciones intracomunitarias. Sustituyendo el viejo modelo de autoliquidación sin contraste previo por un modelo de verificación casi en tiempo real por parte de la propia Administración. Es obligatorio para grandes empresas (facturación superior a 6.010.121,04 euros), inscritos en el REDEME y grupos de entidades, y voluntario para el resto. 

Si LAURÉ diseñó en 1954 un sistema que se autocontrolaba porque cada comprador tenía incentivo a exigir una factura válida para poder deducir, el SII lleva esa misma lógica de verificación cruzada a su versión más avanzada posible: ya no es solo que las partes se controlen entre sí por interés propio, sino que la propia Administración contrasta, casi en el instante, lo que declara el vendedor con lo que declara el comprador, precisamente el mismo horizonte hacia el que la iniciativa europea "IVA en la era digital" (ViDA).

\subsection{Gestión del impuesto}
El sujeto pasivo debe cumplir un conjunto de obligaciones formales que, tomadas en su conjunto, son las que hacen posible que el mecanismo de crédito de factura funcione en la práctica y no solo sobre el papel. 

La primera es la \textbf{autoliquidación periódica}. Es una liquidación trimestral con carácter general y mensual para quienes están inscritos en el Registro de Devolución Mensual (REDEME) o cuentan con un volumen de operaciones superior a 6.010.121,04 euros. Cuando la declaración resulta negativa, el resultado es a compensar en periodos sucesivos (salvo para los inscritos en REDEME), solicitándose la devolución acumulada, en su caso, en la última declaración del ejercicio; existen además regímenes de devolución especiales, como el aplicable a viajeros o a no establecidos.

Junto a la autoliquidación, el sujeto pasivo tiene \textbf{obligaciones censales} (comienzo, modificación y cese de actividad, solicitud y comunicación del NIF), de facturación conforme al Reglamento correspondiente, de contabilidad y libros registro y de presentación de declaraciones asociadas:\footnote{%
    Los libros de registro y contabilidad deben detallar: (i) las facturas emitidas y recibidas; (ii) operaciones de bienes de inversión; y, (iii) de determinadas operaciones intracomunitarias.
} (i) el Modelo 347, declaración recapitulativa de operaciones con terceros que superen, en su conjunto, un determinado umbral anual con cada cliente o proveedor, pensada precisamente para permitir a la Administración cruzar información entre las declaraciones de compradores y vendedores dentro del territorio nacional; (ii) el Modelo 390, declaración-resumen anual; y, (iii) el Modelo 349, declaración de operaciones intracomunitarias, el equivalente doméstico español del sistema europeo de intercambio de información que desarrollaremos en el punto siguiente.

Además del régimen ordinario, existe un catálogo de regímenes especiales, mayoritariamente voluntarios.\footnote{%
    Entre otros: el régimen simplificado, que determina las cuotas devengadas mediante coeficientes en lugar de por diferencia directa entre repercutido y soportado; el régimen especial de agricultura, ganadería y pesca; el de bienes usados, objetos de arte, antigüedades y objetos de colección; el de oro de inversión; el de agencias de viaje; el de recargo de equivalencia, previsto para determinados comerciantes minoristas, cuyos porcentajes exactos desarrollaremos más adelante, junto al resto del catálogo de diferenciación de tipos; el régimen especial del grupo de entidades; el de criterio de caja, que permite retrasar el devengo del IVA repercutido, y el derecho a deducir el soportado, hasta el momento del cobro o pago efectivo, en lugar de al momento ordinario de realización de la operación; y los regímenes de ventanilla única (OSS, antes MOSS), que permiten a un empresario registrarse a efectos de IVA en un único Estado miembro para declarar y liquidar, mediante un único pago, todo el IVA devengado por sus ventas a consumidores finales en cualquier otro Estado de la Unión.
}

\subsubsection{Gestión desde el punto de vista territorial: España, las Comunidades Autónomas y la Unión Europea}

Como ya señalamos, el IVA está cedido a las Comunidades Autónomas en un $50\%$, de acuerdo con el índice de consumo territorial certificado por el INE. El problema, que revela una dificultad técnica genuina de la fiscalidad territorial, es por qué esa cesión no puede articularse de forma directa, atribuyendo a cada Comunidad Autónoma la recaudación de IVA efectivamente generada dentro de su territorio, como sí ocurre, con matices, con otros tributos cedidos. 

La razón es estructural. El IVA se liquida a nivel de empresa, no de transacción individualizada por territorio de consumo. Una gran superficie con sede social en Madrid que vende en toda España no puede, en la práctica, desagregar con precisión en su propia declaración de IVA cuánto de lo repercutido corresponde a un consumidor de Andalucía y cuánto a uno de Cataluña, precisamente porque el propio mecanismo de crédito de factura no exige nunca esa desagregación territorial: solo exige conocer la diferencia entre repercutido y soportado a nivel de sujeto pasivo, no a nivel de punto geográfico de consumo. Por eso España recurre a un mecanismo indirecto y estimativo. En lugar de rastrear la recaudación real generada en cada territorio, distribuye la masa total cedida en proporción a un índice de consumo territorial elaborado por el INE a partir de las estadísticas de gasto de los hogares por Comunidad Autónoma, asumiendo que el consumo declarado por los hogares residentes en cada territorio es la mejor aproximación disponible a dónde se generó, realmente, la carga final del impuesto.\footnote{%
    Es el mismo tipo de proxy estadístico aplicado a la distribución territorial de ingresos, que presentamos al presentar la medición de la regresividad en el debate de POTERBA sobre renta vitalicia ([\authorfont{Ver Tema 4.B.13}]).
}

\begin{specialblock}{\textbf{Límites forales:} ¿Hasta donde alcanza la autonomía fiscal de las Agencias Forales y los Territorios Históricos?}
    Conviene señalar un contraste que revela el peso real de la armonización europea sobre la propia arquitectura fiscal interna española.

    En el País Vasco y en Navarra, bajo sus respectivos regímenes de Concierto y Convenio Económico, las Diputaciones y la Comunidad Foral gestionan y recaudan el IVA generado en su propio territorio, aportando después al Estado una cuota o aportación calculada en proporción a esa recaudación. Evidentemente, un grado de autonomía de gestión notablemente superior al de las Comunidades de régimen común. Pero esa autonomía no se extiende al terreno donde sí resulta decisiva en otros tributos.
    
    Los territorios forales no pueden fijar sus propios tipos de IVA, ni diseñar sus propias exenciones o regímenes especiales al margen de lo armonizado, exactamente lo contrario de lo que ocurre con el IRPF. La razón es la misma que atraviesa todo el bloque. El IVA nació como instrumento de armonización fiscal europea, y esa armonización opera como un límite que ni siquiera la especialidad foral española puede sortear. Es decir, constituye un claro ejemplo de cómo la pertenencia a la Unión Europea puede limitar la autonomía fiscal interna con más fuerza que la propia Constitución.
\end{specialblock}

\subsubsection{Dimensión europea: VIES, Eurofisc y la arquitectura institucional del control transfronterizo}
Como es natural, la cooperación administrativa dentro de la Unión Europea es imprescindible para la gestión de este tributo. De hecho, sin esta el propio régimen transitorio de destino sería, en la práctica, ingobernable. 

El Sistema de Intercambio de Información sobre el IVA (VIES), gestionado a través de un sistema electrónico central que recibe la información transmitida por los sistemas nacionales de cada Estado miembro, permite a cualquier empresa verificar si el número de identificación de IVA de un cliente o proveedor de otro Estado miembro es válido, condición indispensable para poder aplicar correctamente la exención de las entregas intracomunitarias que sostiene todo el régimen de destino. 

Por otro lado, la Red EUROFISC, constituye el mecanismo de intercambio rápido y descentralizado de información específica sobre fraude transfronterizo entre las administraciones tributarias de los Estados miembros, permitiendo detectar patrones sospechosos de operador desaparecido o de carrusel antes de que alcancen la magnitud de episodios anteriores. 

La magnitud del problema que esta arquitectura institucional trata de contener queda reflejada en una estimación conjunta de Europol y el Tribunal de Cuentas Europeo, de gran actualidad: en torno al $2\%$ de los grupos de delincuencia organizada activos en la Unión son responsables de hasta el $80\%$ del fraude intracomunitario de IVA, generando unas pérdidas anuales de recaudación estimadas entre 40.000 y 60.000 millones de euros — una concentración extraordinaria del daño en manos de muy pocos actores, que explica por qué el refuerzo de esta cooperación —incluida, en su desarrollo más reciente, la extensión del acceso a esta información a la propia Fiscalía Europea y a la Oficina Europea de Lucha contra el Fraude (OLAF)— sigue siendo, año tras año, una de las prioridades legislativas más activas de todo el Derecho tributario de la Unión.

\subsection{Esquema básico de liquidación}
Cerramos el primer apartado con un esquema tipo de liquidación: ¿cómo se cuantifica, se liquida y se traslada el IVA a lo largo de una cadena completa, desde el primer eslabón que no soporta ningún impuesto anterior hasta el último que no puede deducir nada de lo que paga?

Para ello, veamos cómo se liquidaría el IVA en una cadena de producción de cuatro eslabones. Por ejemplo, supongamos que se trata del consumo de mobiliario en la cadena sueca IKEA.\footnote{%
    Elegimos deliberadamente una cadena con todos los eslabones plenamente sujetos y no exentos, tributando al tipo general del 21\%, para que quede clara la mecánica sin ninguna de las complicaciones.
}

\textbf{Eslabón 1: Explotación forestal (el inicio de la cadena)}

Una explotación forestal vende un lote de madera a un fabricante de muebles por $1.000 €$. Es el primer eslabón de la cadena, y su rasgo característico es que no tiene nada que restar de lo que repercute.

- Base imponible: 1.000 €.
- IVA repercutido (21\%): 210 €.
- Precio total facturado al fabricante: 1.210 €.
- IVA soportado deducible: 0 €.
- Cuota a ingresar por la explotación forestal: 210 €.

\textbf{Eslabón 2: el fabricante de muebles}

El fabricante compra la madera por 1.000 € + 210 € de IVA (1.210 € en total), soportando así 210 € que podrá deducir íntegramente, al estar afectos por completo a su actividad y no tratarse de ninguno de los bienes con restricción objetiva de deducción que veremos al cierre de este punto. Transforma la madera en una mesa terminada y la vende al mayorista por 2.500 €.

- Base imponible de la venta: 2.500 €.
- IVA repercutido (21\%): 525 €.
- Precio total facturado al mayorista: 3.025 €.
- IVA soportado deducible (el pagado al comprar la madera): 210 €.
- Cuota a ingresar por el fabricante: 525 − 210 = 315 €.

Fijemonos en que los 315 € que el fabricante ingresa a Hacienda son, exactamente, el 21\% de su propio valor añadido: la diferencia entre el precio de venta (2.500 €) y el precio de compra de su materia prima (1.000 €), es decir, el 21\% de 1.500 € = 315 €. Comprobamos pues que el crédito de factura hace irrelevante el número de fases de la cadena: cada eslabón, sin excepción, termina ingresando exactamente el 21\% de lo que él mismo ha añadido de valor, ni un euro más.

\textbf{Eslabón 3: el mayorista}

El mayorista compra la mesa por 2.500 € + 525 € de IVA (3.025 € en total), soportando 525 € deducibles. La revende a una tienda de muebles minorista por 3.200 €.

- Base imponible: 3.200 €.
- IVA repercutido (21%): 672 €.
- Precio total facturado a la tienda: 3.872 €.
- IVA soportado deducible: 525 €.
- Cuota a ingresar por el mayorista: 672 − 525 = 147 € — de nuevo, exactamente el 21% de su valor añadido (3.200 − 2.500 = 700 €; 21% de 700 = 147 €).

#### Eslabón 4: la tienda minorista (el penúltimo eslabón)

La tienda compra la mesa por 3.200 € + 672 € de IVA (3.872 € en total), soportando 672 € deducibles. La vende, finalmente, a un particular por 4.000 €.

- Base imponible: 4.000 €.
- IVA repercutido (21%): 840 €.
- Precio total que paga el consumidor: 4.840 €.
- IVA soportado deducible: 672 €.
- Cuota a ingresar por la tienda: 840 − 672 = 168 € — el 21% de su valor añadido (4.000 − 3.200 = 800 €; 21% de 800 = 168 €).

Aquí está el punto donde la cadena se detiene de forma definitiva. El particular que compra la mesa por 4.840 € no es empresario ni profesional, no repercute nada a nadie después de él, y no tiene ningún derecho de deducción. Los 840 € de IVA que ha pagado dentro de esos 4.840 € son, para él, un coste definitivo y sin retorno. Propiedad que distingue al consumidor final de cualquiera de los cuatro eslabones empresariales anteriores, cada uno de los cuales ha manejado en sus libros cifras de IVA mucho mayores pero ha terminado ingresando una cantidad mucho menor (neta).

Si sumamos ahora las cuatro cuotas netas efectivamente ingresadas a la Hacienda Pública a lo largo de toda la cadena: 210 € + 315 € + 147 € + 168 € = 840 €. Probamos que es, exactamente, el IVA que pagó el consumidor final en su única transacción: 840 €. Ni más, ni menos, con independencia de que la mesa haya pasado por cuatro manos distintas antes de llegar al salón de quien finalmente la disfruta.

Conviene hacer dos precisiones que completan el esquema y que, a diferencia de la prorrata, no dependen de que existan exenciones. Primera, la ley excluye de la deducción, con independencia de la actividad del sujeto pasivo, las cuotas soportadas por la adquisición de joyas, alhajas, piedras preciosas, alimentos, bebidas, tabaco, y espectáculos y servicios de carácter recreativo, una restricción que no responde a ninguna exención de la operación en sí, sino a una presunción legal de que estos bienes concretos, aun comprados formalmente por una empresa, tienden a servir a un consumo de naturaleza mixta o personal difícil de deslindar del empresarial, y que el legislador ha preferido excluir de raíz antes que fiscalizar caso por caso. Segunda, en nuestro ejemplo, todos los eslabones tenían más IVA repercutido que soportado; pero esto no tiene porque ser así. Imaginemos que el fabricante de muebles, en el mismo periodo en que compró la madera, hubiera invertido además en maquinaria nueva por 10.000 € más IVA (2.100 €), sin haber llegado todavía a vender ninguna mesa. Su IVA soportado (210 € de la madera + 2.100 € de la maquinaria = 2.310 €) superaría entonces a su IVA repercutido (0 €, si no ha vendido nada aún), generando un saldo a su favor. Ese saldo no se pierde, se compensa en las declaraciones-liquidaciones de periodos sucesivos hasta agotarse, o, si al cierre del ejercicio sigue existiendo saldo pendiente, se solicita su devolución. 

El esquema de liquidación, por tanto, no es solo una secuencia de pagos netos crecientes. Cualquier empresa que invierta intensamente antes de facturar (nueva creación) puede situarse durante varios periodos en posición acreedora frente a la Hacienda Pública, sin que ello altere en absoluto la propiedad de neutralidad.







\clearpage
\section{Impuesto sobre el Valor Añadido: Especificidades}
\puntosclave{
    
}

Como podemos intuir, el esquema de liquidación que acabamos de cerrar es una ficción útil. En la práctica, ningún operador económico español tributa exactamente así está sembrado de excepciones a esa regla general. 

No obstante, el régimen limpio presentado hasta ahora sigue siendo absolutamente necesario como punto de partida. Se trata del benchmark frente al cual medimos cada una de las divergencias que vamos a pasar a analizar. Recorreremos cuatro preguntas: (i) por qué existen estas divergencias, y qué las justifica; (ii) qué es, técnicamente, una exención y cómo se gestiona en la práctica española; (iiii) qué es la diferenciación de tipos y cómo se gestiona; y, (iv) un balance general que enfrenta ambos instrumentos entre sí y los sitúa en el debate de actualidad más reciente: Libro Blanco de 2022.

\subsection{¿Por qué existen estas divergencias?}
Las exenciones y los tipos reducidos no responden al mismo tipo de argumento, y tratarlos como si fueran dos variantes de una misma lógica es un error común que debemos evitar. 

\subsubsection{Exenciones: Interpretación esctrica}
Por un lado, las exenciones (art. 20 LIVA) se agrupan, casi todas, en dos familias con fundamentos completamente distintos entre sí: (i) imposibilidad técnica de medición, el caso paradigmático de los servicios financieros; y, (ii) de interés general o social, como la sanidad, educación, asistencia social.

La cuestión metodológica es muy fácil de comprender cuando analizamos el caso de los servicios financieros. Como es natural, el margen implícito entre el tipo de interés pagado y el cobrado, resulta imposible de descomponer transacción por transacción en la práctica. Respecto a la segunda familiar, la Directiva 2006/112/CE del IVA (art. 132; art. 20 LIVA) agrupa bajo la rúbrica exenciones aplicables a ciertas actividades de interés general la sanidad, la educación, los servicios sociales y culturales, y las actividades sin ánimo de lucro de organizaciones políticas, sindicales, religiosas o filantrópicas. Naturalmente, la lógica de fondo es la de los bienes de mérito o bienes preferentes (MUSGRAVE), pues el legislador europeo considera que la sanidad y la educación deben quedar al margen del cálculo puramente mercantil del precio final precisamente porque su consumo se considera socialmente deseable con independencia de la capacidad de pago de quien lo recibe.

Respecto a la segunda cuestión, de mucha menor claridad, el TJUE ha construido, desde finales de los años ochenta (C-348/87, 1989), una doctrina consolidada según la cual las exenciones del IVA deben interpretarse de forma estricta. Como determina el tribunal, constituyen excepciones al principio general según el cual el IVA se percibe sobre toda entrega de bienes y toda prestación de servicios efectuada a título oneroso por un sujeto pasivo.

Esto significa algo con consecuencias prácticas muy concretas. Ante cualquier duda razonable sobre si una operación concreta encaja o no dentro del perímetro de una exención, la interpretación correcta es la que menos amplía el ámbito de la excepción. La jurisprudencia reciente del TJUE ofrece ejemplos muy concretos de esta lógica en acción: una sentencia de 2020 (asunto C-513/20) tuvo que precisar que la exención de curas termales solo se aplica cuando media prescripción médica; otra, de 2020 también (C-48/19), exigió que el asesoramiento telefónico sobre salud persiga una finalidad terapéutica concreta para beneficiarse de la exención sanitaria; y, otra de 2019 (C-42/18) acotó con precisión el alcance exacto de la exención de operaciones financieras frente a servicios accesorios que, aunque prestados por una entidad financiera, no encajan estrictamente en el perímetro exento.

El mensaje que esta jurisprudencia transmite es que la propia arquitectura jurídica europea trata las exenciones como lo que son: una desviación deliberada y acotada de la norma general, nunca un régimen alternativo de igual jerarquía. 

\subsubsection{Diferenciación de tipos: el argumento distributivo, y su límite comunitario —la lista cerrada del Anexo III}
Por otro lado, los tipos reducidos responden casi siempre a la preocupación distributiva sobre bienes de consumo generalizado, terreno donde entra en juego la tensión ATKINSON-STIGLITZ/CORLETT-HAGUE que ya tratamos en otros temas ([\authorfont{Ver Tema 4.B.13}]). De esta forma, el fundamento de los tipos reducidos es de naturaleza muy distinta: no corrige ningún problema de medición ni persigue eximir por interés general, sino que busca modular la carga fiscal sobre bienes de consumo generalizado por razones distributivas.

Lo que conviene fijar aquí es que esta facultad de los Estados miembros no es ilimitada. La Directiva 2006/112/CE contiene una lista cerrada y tasada de categorías de bienes y servicios sobre los que los Estados miembros pueden, si lo desean, aplicar un tipo reducido (alimentos, productos farmacéuticos, libros, transporte de viajeros, entre otras), sin que ningún Estado pueda inventar libremente nuevas categorías fuera de esa lista.\footnote{%
    El episodio que mejor ilustra hasta qué punto esta lista cerrada puede generar resultados que hoy resultan casi absurdos es el de los libros electrónicos. La Gran Sala del TJUE (asunto C-390/15), tuvo que resolver si Polonia podía aplicar a los libros electrónicos el mismo tipo reducido que ya aplicaba a los libros en papel. La respuesta, sorprendentemente, fue que no. La razón es que, siguiendo estrictamente la letra del Anexo III vigente en ese momento, la norma hablaba de libros en cualquier tipo de soporte físico; por tanto el Tribunal consideró que un libro electrónico, al suministrarse por vía electrónica, quedaba excluido de la categoría admitida, con independencia de que su contenido fuera exactamente el mismo.

El resultado resultaba tan difícil de justificar que la propia Unión reaccionó reformando el Anexo III apenas un año después, mediante la Directiva (UE) 2018/1713, para permitir expresamente la aplicación de tipos reducidos también a las publicaciones electrónicas.
}

Como es evidente, la justificación a esta restricción es de neutralidad. Si cada Estado miembro pudiera aplicar tipos reducidos a cualquier bien que decidiera unilateralmente, la propia armonización fiscal quedaría vaciada de contenido, reabriendo la puerta a distorsiones de competencia entre Estados miembros que la propia existencia del IVA armonizado pretendía cerrar

\subsection{¿Qué es una exención y cómo se trata?}
La cuestión fundamental de las exenciones es su distinción técnica. Una operación exenta no repercute IVA, pero esa exención puede ser: plena, si permite, además, deducir el IVA soportado en las adquisiciones relacionadas con la actividad exenta; o, limitada, si, junto con la no repercusión, arrastra la pérdida del derecho a deducir.\footnote{%
    \textbf{NOTA.-} No debemos confundir exención con la \textbf{no sujeción}. El resultado práctico puede verse idéntico a primera vista: en la no sujeción, el supuesto de hecho que activaría la obligación tributaria ni siquiera llega a producirse; y, en la exención, la obligación ha nacido, pero la ley neutraliza sus efectos por una razón específica y acotada. Sin embargo, la diferencia no es solo semántica: solo la exención está sujeta a la doctrina de interpretación estricta del TJUE, precisamente porque solo ella constituye, en sentido técnico, una excepción a una regla que de otro modo se aplicaría.
}

La legislación reserva las exenciones plenas, casi en exclusiva, a las exportaciones y operaciones asimiladas, a las entregas intracomunitarias, y a las operaciones relacionadas con zonas y depósitos francos. Todas las demás exenciones del catálogo español son limitadas, y es precisamente esa limitación, con su pérdida asociada del derecho a deducir, la que genera las consecuencias que nos ocupan.

\subsubsection{Catálogo: Exenciones por finalidad}
El art. 20 LIVA agrupa las exenciones interiores en varios bloques por finalidad.

\paragraph{Bienes preferentes: Interés general}
Pertenece, sin excepción, a la familia del interés general: no hay aquí ningún problema de medición del margen, solo una decisión deliberada de excluir del circuito de IVA el consumo de bienes considerados socialmente preferentes.
\begin{lnum}
    \item \textbf{Ámbito sanitario y educativo}. Servicios de hospitalización y asistencia sanitaria, diagnóstico, prevención y tratamiento de enfermedades; transporte de heridos y enfermos en ambulancia; entregas de sangre, plasma y elementos del cuerpo humano con fines médicos o de investigación; servicios odontológicos y de prótesis dental; educación de la infancia y juventud en todos sus niveles, formación y reciclaje profesional, y clases particulares sobre materias del plan de estudios oficial impartidas por personas físicas;
    \item \textbf{Servicios sociales, culturales y deportivos}. Asistencia social prestada por entidades de derecho público o establecimientos privados de carácter social (tercera edad, minorías étnicas, exreclusos); prestaciones a los propios miembros de organizaciones sin finalidad lucrativa, sin más contraprestación que las cuotas estatutarias; bibliotecas, museos y representaciones teatrales de entidades públicas o de carácter social; servicios a deportistas. 
\end{lnum}


\paragraph{Fundamento técnico: Problemas de medición}
Recuerda que el fundamento aquí no es social, sino puramente técnico —la dificultad, ya trabajada con el Fondo Monetario Internacional y su propuesta de Financial Activities Tax, de descomponer el margen de intermediación financiera transacción por transacción—.
\begin{lnum}
    \item \textbf{Operaciones de seguro y financieras}. Aquí es donde entra la exención por imposibilidad de medición que ya desarrollamos con detalle en el Tema 13: seguro, reaseguro y capitalización; constitución y transmisión de préstamos y créditos; operaciones con depósitos en efectivo; transmisión de valores; y operaciones con divisas, billetes y monedas de curso legal;
    \item 
\end{lnum}



Exenciones inmobiliarias. Entregas de terrenos rústicos y no edificables; entregas de terrenos como aportación a Juntas de Compensación; segundas y ulteriores entregas de edificaciones tras su terminación; y arrendamientos y constitución de derechos reales de goce destinados a vivienda. Esta familia merece desarrollo propio, porque es la que mejor ilustra —con datos y mecanismos reales, no solo con teoría— todo lo que vamos a construir en el epígrafe ii.5.

Otras exenciones. Servicios postales; prestaciones entre miembros de uniones y agrupaciones autónomas (incluidas las Agrupaciones de Interés Económico); loterías, apuestas y juegos gestionados por SELAE, ONCE y los organismos autonómicos correspondientes.

Exenciones técnicas. Entregas de bienes cuya adquisición o afectación hubiera determinado, en su día, la exclusión total del derecho a deducir para el transmitente. [Probable] Esta última familia merece una mención especial porque no encaja del todo en ninguna de las dos categorías de fundamento que fijamos. No persigue ningún interés general ni responde a ningún problema de medición, sino que evita, con pura lógica interna del sistema, que un bien que nunca generó derecho a deducción en su adquisición vuelva a tributar íntegramente al revenderse — un ajuste de coherencia interna del propio mecanismo, más que una divergencia sustantiva sobre qué merece protegerse.

\subsubsection{Implicaciones: ¿cómo afecta esta figura?}
Como sabemos, la deducción exige que los bienes y servicios adquiridos se destinen a operaciones sujetas y no exentas. Por tanto, cuando una parte de la actividad de un sujeto pasivo está sujeta y no exenta, y otra parte está exenta de forma limitada, surge un problema que el esquema anterior no contemplaba en absoluto. ¿Qué proporción del IVA soportado en las adquisiciones comunes a ambas actividades puede deducirse? 

\paragraph{Implicaciones prácticas: ¿cómo se ajustan las operaciones en un mismo sujeto pasivo?}
La respuesta es la regla de prorrata. Es el ajuste necesario que surge como consecuencia de la posibilidad de que convivan, en un mismo sujeto pasivo, actividades con derecho a deducir y actividades sin ese derecho. La Ley distingue dos modalidades:
\begin{la}
    \item \textbf{General}. La prorrata general aplica un único porcentaje a todo el IVA soportado común, con independencia de a qué actividad concreta se destine cada adquisición. Este porcentaje se calcula como el cociente entre el importe anual de operaciones con derecho a deducir y el volumen total de operaciones del sujeto pasivo; y,
    \item \textbf{Especial}. La prorrata especial, que exige deducir según el destino real de cada adquisición. El sujeto tendrá, por tanto, dos opciones: (i) deducir íntegramente si se destina en exclusiva a la actividad con derecho a deducción; o, (ii) nada si se destina en exclusiva a la actividad exenta. Y, además, solo el porcentaje de prorrata general para lo que se destine a ambas actividades a la vez. En general, la prorrata especial es optativa. Sin embargo, se convierte en obligatoria cuando la prorrata general arroje una deducción que exceda en más de un $20\%$ a la que resultaría de aplicar la especial, precisamente para evitar que el sujeto pasivo se beneficie, mediante la elección del método menos preciso, de una deducción sistemáticamente superior a la que su actividad real justificaría.
\end{la}

Si recuperamos el anterior, podemos ver como cambia en caso de que alguna transacción esté sujeta a una exención.

Por ejemplo, supongamos que un empresario arrienda, simultáneamente, locales comerciales (sujeta) y viviendas (exenta sin derecho a deducción). Este empresario factura $4.000.000 €$ por los locales y $1.000.000 €$ por las viviendas.

- Prorrata: operaciones con derecho a deducción / total = 4.000.000 / 5.000.000 = 80%.
- Adquisiciones de bienes y servicios al tipo general (21%), por un total de 3.000.000 €, distribuidas así: 1.000.000 € afectos en exclusiva a los locales, 1.000.000 € afectos en exclusiva a las viviendas, y 1.000.000 € de gastos comunes a ambas actividades (por ejemplo, servicios de administración de fincas que atienden simultáneamente a todo el edificio).
- IVA repercutido: 4.000.000 × 21% = 840.000 €.
- IVA soportado total: 3.000.000 × 21% = 630.000 €.
- Bajo prorrata general: 630.000 × 80% = 504.000 € deducibles.
- Bajo prorrata especial: 210.000 € (los afectos a locales, deducibles al 100%) + 0 € (los afectos a viviendas, deducibles al 0%) + 168.000 € (los comunes, 210.000 × 80%) = 378.000 € deducibles.
- La diferencia entre ambos métodos —504.000 € frente a 378.000 €— representa un 33,3% de exceso de la prorrata general sobre la especial, muy por encima del umbral del 20%: la prorrata especial es obligatoria.
- IVA a ingresar = 840.000 − 378.000 = 462.000 €.

Comparémoslo con lo que habría ocurrido si toda la actividad hubiera estado sujeta y no exenta. El empresario habría deducido el $100\%$ de los $630.000€$ soportados, ingresando solo $210.000€$. Por tanto, la existencia de la exención sobre las viviendas ha encarecido en $252.000€$ la cuota a ingresar de este empresario, sin que ese incremento corresponda a un aumento real de su actividad económica. Se trata, en esencia, del coste de que una parte de su negocio quede fuera del circuito de neutralidad que LAURÉ diseñó.

\paragraph{Implicaciones teóricas: Cascada en miniatura}
Como vemos, la exención reintroduce en miniatura la acumulación piramidal que el crédito de factura eliminó. Para demostrarlo, veamos un mecanismo real. El caso más ilustrativo es el de las segundas entregas de edificaciones.

Imagina un promotor que rehabilita un edificio de oficinas, soportando $200.000€$ de IVA en obras y materiales, y lo vende después a una empresa que va a destinarlo, íntegramente, a su propia actividad sujeta y no exenta. Si la venta del edificio queda exenta (regla general para segundas entregas), el promotor no puede deducir esos $200.000€$ de IVA que soportó en la rehabilitación, y los repercute, de forma oculta y sin desglose visible, dentro del precio final del inmueble. La empresa compradora, al no recibir ninguna factura con IVA desglosado, tampoco puede deducir nada de esos $200.000€$ escondidos en el precio, aunque vaya a destinar el edificio en exclusiva a una actividad con pleno derecho a deducción. 

Por tanto, el IVA no deducido queda enterrado dentro del precio, y se traslada a la siguiente fase de la cadena. De nuevo, la misma acumulación que LAURÉ diseñó su sistema para eliminar.

Como es normal, la LIVA no ignora este problema y permite (art. 20), bajo determinadas condiciones, la renuncia a la exención.\footnote{
Las condiciones son que el adquirente sea empresario o profesional con derecho a deducción total o parcial en función del destino previsible del inmueble, y que el transmitente lo comunique fehacientemente al adquirente con carácter previo o simultáneo a la entrega.
} La lógica económica de esta renuncia es exactamente la que el ejemplo anterior predice. Al promotor y al comprador les conviene conjuntamente renunciar a la exención y tributar por IVA de forma ordinaria, porque así el promotor recupera sus $200.000€$ de IVA soportado en la rehabilitación, y el comprador puede deducir el IVA que su vendedor le repercuta, si va a destinar el inmueble a una actividad con derecho a deducción. La exención, pensada en su origen para proteger determinadas transacciones inmobiliarias (vivienda), se convierte en un obstáculo puramente contraproducente cuando ambas partes de la operación son empresarios con plena capacidad de moverse dentro del circuito de neutralidad del IVA.

Si la exención fuera siempre beneficiosa para el sujeto pasivo, nadie renunciaría a ella. No obstante, el hecho de que la propia norma prevea un procedimiento para renunciar a un beneficio fiscal confirma que la exención limitada no es, en absoluto, un regalo sin coste: es un instrumento cuyo efecto neto depende de la posición de cada parte dentro de la cadena, y que puede resultar tanto favorable como directamente perjudicial frente al régimen general plenamente neutro, que presentamos antes.

\subsection{¿Qué es la diferenciación de tipos y cómo se trata?}
Antes de recorrer el catálogo, hay que señalar algo que separa radicalmente este instrumento de la exención, porque es fácil confundirlos bajo la etiqueta común de beneficio fiscal: \textbf{aplicar un tipo reducido a una operación no rompe la cadena de deducción}. 

Una empresa que vende pan al $4\%$ sigue siendo sujeto pasivo plenamente integrado en el circuito de neutralidad. Repercute IVA, aunque a un tipo menor, y deduce con normalidad todo el IVA que soporta en sus compras, con independencia de a qué tipo tributen esas compras. Por tanto, no hay aquí ningún equivalente a la prorrata, ni ninguna cascada en miniatura; lo único que cambia es la cuantía de lo repercutido, no la mecánica de la deducción.

\subsubsection{Catálogo: General, reducido y superreducido}

Recordarás, del Bloque 1, los tres tipos que ya anticipamos sin desarrollar: general (21\%), reducido (10\%) y superreducido (4\%). El reducido se aplica, entre otros, a la transmisión de viviendas, a determinados productos vegetales, a los servicios de hostelería y restauración, al transporte de viajeros, y a los alimentos para consumo humano y animal no incluidos en el superreducido. El superreducido se aplica a alimentos básicos, medicamentos de uso humano, libros y prensa sin contenido publicitario dominante, material escolar, vivienda de protección oficial, prótesis y órtesis, y determinados servicios de atención a la dependencia.

A este catálogo se han sumado, desde 2023, dos nuevos tipos: (i) del $0\%$, sobre alimentos básicos; y, (ii) un tipo del $5\%$ (rebajado desde el $10\%$), para pastas y aceites. Hasta entonces inexistentes, se implementaron como medida de contención de la inflación desatada por la invasión rusa de Ucrania de forma transitoria y, como es natural en este tipo de cuestiones, se fueron prorrogando sucesivamente durante 2023 y 2024. El propio Gobierno anunció, en junio de 2024, el inicio de una retirada gradual y escalonada de esta rebaja.\footnote{
\textcolor{red}{\textbf{OJO ->}} Conviene verificar si esta senda de retirada gradual se ha completado tal y como estaba prevista o si ha vuelto a prorrogarse.
}

\subsubsection{Implicaciones: ¿cómo altera el esquema de liquidación?}
Como decíamos, aunque la diferenciación de tipos no rompe la cadena de deducción, sí altera de forma muy visible la posición de tesorería de determinados sujetos pasivos. 

\paragraph{Implicaciones prácticas: posición estructuralmente acreedora}
Imagina una panadería que vende pan al tipo actualmente vigente para ese producto, mientras que buena parte de sus costes tributan al tipo general del $21\%$. Si en un periodo determinado la panadería factura $50.000€$ en ventas de pan ($4\%$) y soporta $15.000€$ de gastos generales al $21\%$ ($3.150€$ de IVA soportado), su IVA repercutido asciende a solo $2.000€$.  La panadería se encuentra, en ese periodo, con más IVA soportado que repercutido. Una posición estructuralmente acreedora frente a la Hacienda Pública, que arrastrará mediante compensación en periodos sucesivos o solicitará en devolución al cierre del ejercicio. 

Esta es la consecuencia estructural de la diferenciación de tipos. Cualquier sector cuyo producto final tribute a un tipo sensiblemente inferior al de sus propios costes de producción se convierte, de forma sistemática, en un generador recurrente de saldos a devolver, con la carga administrativa y el coste de tesorería que ello supone. Un coste real del sistema que rara vez se cuantifica cuando se debate, en términos puramente distributivos, la conveniencia de mantener o ampliar un tipo reducido.

\paragraph{Implicaciones teóricas: }
Además, recordando las aportaciones de ATKINSON y STIGLITZ y de CORLETT y HAGUE, podemos encontrar una serie de implicaciones examinando el catálogo español.

Por ejemplo, los alimentos básicos y medicamentos al tipo superreducido encajan en la lógica ATKINSON-STIGLITZ de equidad pura: la preocupación distributiva convencional sobre bienes de primera necesidad cuyo consumo pesa proporcionalmente más sobre la renta de los hogares menos favorecidos. Por otro lado, los servicios de hostelería y restauración encajan en la lógica CORLETT-HAGUE de sustitución del ocio: comer fuera de casa libera tiempo que, de otro modo, se destinaría a cocinar, permitiendo a quien lo consume dedicar más horas al trabajo remunerado (\textit{services à la personne}). El transporte de viajeros admite la misma lectura, reduce el tiempo de desplazamiento y facilita, de forma indirecta, la disponibilidad de tiempo laboral.

Los libros, la prensa y el material escolar, situados en el tipo superreducido, son más difíciles de encajar limpiamente en cualquiera de las dos lógicas. No son, en sentido estricto, bienes de supervivencia material como los alimentos, pero tampoco son claramente sustitutivos del ocio en el sentido de liberar tiempo de trabajo, más bien lo contrario. Su tipo reducido responde, con mayor probabilidad, a un tercer argumento: una consideración de externalidad cultural positiva, próxima al argumento de bien de mérito que ya trabajamos con MUSGRAVE ([\authorfont{Ver Tema 4.B.13}]).

\paragraph{Implicaciones administrativas: Diferenciación en frontera}
Por último, conviene señalar que cualquier catálogo de tipos diferenciados genera, inevitablemente, problemas en la delimitación de los bienes en frontera. 

El propio catálogo de alimentos básicos que acabamos de recorrer ofrece un ejemplo perfecto y actual. Mientras que las frutas, verduras y legumbres frescas se benefician del tipo más reducido, las conservas de esos mismos productos, el arroz precocido, o las patatas prefritas congeladas, quedan excluidas de ese beneficio y tributan al 10\%. Este tipo de fronteras, generan un volumen constante de consultas vinculantes a la Dirección General de Tributos y de litigios ante los tribunales económico-administrativos, dedicados a resolver preguntas tan aparentemente triviales como si un determinado producto alimenticio conserva o no la condición de producto natural a efectos del Código Alimentario que la propia norma del IVA invoca como criterio de remisión. 

Esto genera inevitablemente costes administrativos, litigiosidad e inseguridad jurídica en el margen, que deben ser correctamente ponderados a la hora de valor su adecuación y llevar a cabo un análisis coste-beneficio. Los costes ocultos aparecerán en cualquier sistema que renuncie a la uniformidad; y, además, cuantas más categorías distintas se crean, más fronteras hay que vigilar, y ninguna frontera deja jamás de ser conflictiva.


\subsection{Balance general: ¿cuál de los dos instrumentos es preferible y qué es lo que observamos?}
Con todo lo expuesto, y antes de entrar en debates de actualidad, conviene cerrar una pregunta. Entre exención y tipo reducido, ¿cuál es el instrumento técnicamente superior para perseguir un mismo objetivo de política social? 

La respuesta es clara y no admite ambigüedad. Si el objetivo del IVA es que un bien o servicio llegue al consumidor final con la menor carga fiscal posible sin distorsionar las decisiones de producción, un tipo reducido es siempre superior a una exención, porque preserva el derecho a deducir del proveedor y evita la cascada en miniatura.

Por tanto, debemos preguntarnos por qué España utiliza con tanta frecuencia la exención. Evidentemente, no responde a ninguna preferencia deliberada por el instrumento inferior; sino a que el propio marco jurídico europeo no ofrece esa opción para estas categorías. Las exenciones del art. 132 de la Directiva (interés general) y los tipos reducidos del Anexo III (lista cerrada) son dos regímenes jurídicos distintos, y el tipo cero generalizado para categorías nuevas está cerrado como opción dentro del sistema armonizado vigente. Por tanto, España no eligió el peor instrumento: la Unión Europea no le ofrece el mejor para este terreno concreto.

Ahora bien, la partida no está ganada. De acuerdo con uno de los Informes más recientes sobre tributación indirecta en el ámbito de la Unión, el uso relativo de ambos instrumentos en España frente al resto de la Unión es muy ilustrativo: el tipo implícito del IVA español es de los más bajos de toda la Unión Europea, superior únicamente al de Rumanía e Italia. El tipo implícito de IVA se calcula con la recaudación efectiva expresada como proporción del consumo total, por tanto es la medida que mejor capta el efecto agregado de exenciones y tipos reducidos combinados.

\subsubsection{Propuestas de reforma: Argumentos a favor de la convergencia}
¿Qué revelan estos datos? En esencia, que España no es un caso marginal dentro del uso de estas divergencias, sino uno de los países que más las ha desarrollado. Con un catálogo de tipos reducidos y superreducidos particularmente extenso, España se sitúa sistemáticamente por debajo de lo que su tipo general nominal del $21\%$ haría esperar. 

De allí las frcuentes alusiones al caso Español en los debates europeos relativos a la armonización de la fiscalidad indirecta. Por tanto, ¿qué debemos hacer para asegurar la convergencia?

\paragraph{Libro Blanco de 2022: ¿qué debemos hacer para converger con la Unión?}
El Comité de Personas Expertas para la Reforma del Sistema Tributario (2021) dentro del Componente 28 del Plan de Recuperación, Transformación y Resiliencia, entregó su Libro Blanco tras diez meses de trabajo.\footnote{%
    Presidido por CARBONELL, catedrático emérito de Economía Aplicada de la Universidad Rey Juan Carlos, el Comité reunió a dieciséis vocales de reconocida trayectoria en Economía y Derecho Financiero y Tributario (LABORDA, RUIZ-ALMENDRAL, CASASNOVAS, LAGO PEÑAS y LABANDEIRA), produjo un informe de cerca de $700$ páginas. El propio Ruiz-Huerta insistió, en la rueda de prensa de presentación, en que el documento no es una reforma fiscal, sino un conjunto de propuestas y sugerencias cuyas decisiones últimas pertenecen al ámbito político.
} En esencia, el Informe articula una propuesta de reforma con $37$ recomendaciones generales y $118$ propuestas concretas de cambio. 

En lo relativo al IVA, el Comité parte del dato de tipo implícito bajo y lo vincula explícitamente a la existencia de gastos tributarios ineficientes e inequitativos que sería necesario reducir para cerrar el desfase con la UE en la recaudación de los grandes tributos, como ocurre especialmente en el IVA. Por ello, no es de extrañar que la propuesta central sea, precisamente, eliminar progresivamente los tipos reducido y superreducido, y converger hacia un tipo único. 

El propio Comité ofrece dos escenarios alternativos según cuál sea la prioridad política: (i) si se mantiene el tipo general en el $21\%$ y se eliminan los reducidos, la recaudación adicional se estima en $27.100$ millones de euros; si, en cambio, (ii) se prioriza la neutralidad recaudatoria, el tipo general podría reducirse hasta aproximadamente el $15\%$ manteniendo la recaudación constante. Una cifra que, de aplicarse, situaría a España muy por debajo del tipo general medio de la Unión Europea, invirtiendo por completo el diagnóstico de partida sobre el bajo tipo implícito español.

\paragraph{Comisión LAGARES (2015): El precedente directo, un largo recorrido sin implementación}
Conviene señalar que el Libro Blanco de 2022 no fue la primera vez que un comité de expertos español recomendó esta misma dirección. Ya en 2014, la llamada Comisión LAGARES había planteado una reforma con el mismo diagnóstico de fondo.\footnote{%
    LAGARES llegó a calificar la situación entonces vigente con una frase que resume por sí sola el problema: más del $40\%$ de los bienes están fuera de esa franja del $21\%$.
} 

En esta ocasión, la Comisión propuso que solo el turismo, la vivienda y el transporte de personas y mercancías conservaran el tipo reducido, con el resto de bienes hoy beneficiados por tipos del $4\%$ o el $10\%$ saltando directamente al tipo general. La subida estimada por la Comisión LAGARES, mucho más modesta (entre $5.000$ y $6.000$ millones), tampoco llegó a implementarse en aquel momento. 

Como vemos, dos comités de expertos han llegado a la misma conclusión de fondo: el catálogo español de tipos reducidos es excesivamente amplio y debería recortarse. Sin embargo, ambos han visto sus recomendaciones archivadas por el poder político sin traducción legislativa sustancial. Así que, más que aceptarlo como un \textit{fallo} político, lo que cabría preguntarse es por qué.

\subsubsection{Preservación del \textit{status quo}: Argumentos en contra de la convergencia}
La crítica de fondo a la propuesta del Comité de 2022, y la anterior por extensión, no proviene de un rechazo a su diagnóstico técnico, sino de su impacto distributivo.

\paragraph{Efectos distributivos: Inequidad y mecanismos compensatorios}
Eliminar tipos reducidos sobre bienes de primera necesidad encarece, en términos proporcionales, mucho más a los hogares de renta baja que a los de renta alta. 

Es por ello que el propio Comité, consciente de esta objeción, la incorporó como condición explícita de su propia propuesta y no como un añadido posterior. Cualquier ampliación de la base del IVA debería ir acompañada de \textbf{mecanismos de compensación} específicos para los hogares de menor renta; ya sea mediante prestaciones directas, o mediante una deducción reembolsable en el IRPF.\footnote{%
    Conectando así la reforma del IVA con la propia reforma del IRPF que el Comité proponía en paralelo, e incluyendo un crédito fiscal por rendimientos del trabajo de tipo EITC que ya vimos al hablar de los impuestos negativos sobre la renta ([\authorfont{Ver Tema 4.B.16}]).
} 

\paragraph{Riesgos políticos: Fragmentación y descordinación}
Por otro lado, y como crítica más objetivo que sustantiva, condicionar la subida del IVA a una compensación simultánea y bien calibrada exige una coordinación legislativa que ningún proceso parlamentario español reciente ha demostrado ser capaz de sostener con la disciplina necesaria. El riesgo, señalado por buena parte de los análisis, es que la subida del IVA (inmediata y automática) llegue a implementarse mientras la compensación prometida (sujeta a negociación presupuestaria) se demore o diluya, dejando a los hogares de renta baja expuestos al primer efecto sin haber recibido todavía el segundo. En el fondo, es el mismo riesgo de descoordinación entre teoría de diseño óptimo y economía política de la implementación que atraviesa buena parte de los grandes debates de reforma fiscal: \textbf{el diagnóstico técnico rara vez es el obstáculo; la secuenciación política casi siempre lo es.}




\clearpage
\section{Impuestos especiales}
\puntosclave{
    
}


Cerramos el tema con el tercer gran pilar de la imposición indirecta española, y el que exige el mayor cuidado de los tres, porque conviven en él seis figuras normativas distintas —cuatro dentro del alcohol, más hidrocarburos, tabaco, electricidad, carbón y medios de transporte— cada una con su propia lógica y su propio debate. Recorremos este bloque en tres tiempos: primero, la justificación teórica de por qué existe siquiera una imposición selectiva distinta y complementaria del IVA general (este primer apartado); después, cada una de las principales figuras impositivas con el detalle que merece; y, para cerrar, una valoración conjunta que las compare entre sí y las sitúe dentro del sistema completo de imposición indirecta que hemos construido a lo largo de todo el tema.

## 1. Introducción y justificación teórica de la imposición especial

### 1.1. Por qué no basta con el IVA: la doble función recaudatoria y correctora

[ICEX-CECO] El gravamen sobre el consumo no puede quedar confiado en exclusiva a un impuesto general como el IVA. Es necesaria la existencia de tributos que graven selectivamente el consumo de bienes específicos, y esa necesidad responde a una doble función: recaudatoria, por un lado —los Impuestos Especiales aportan una cifra muy sustancial de ingresos públicos, con hidrocarburos y tabaco como las dos partidas mayores—, y de instrumento de política económica, por otro, en la medida en que persiguen reducir las externalidades negativas asociadas al uso o producción de determinados bienes.

### 1.2. El fundamento pigouviano, desarrollado con el rigor que merece

Ya trabajamos, en el Tema 13, el fundamento clásico de cualquier impuesto correctivo: Arthur Cecil Pigou mostró que, cuando el consumo de un bien genera costes sociales que exceden el coste privado que asume quien lo consume, el mercado produce y consume ese bien en una cantidad superior a la socialmente óptima, y un impuesto igual al coste externo marginal restaura la eficiencia obligando al consumidor a internalizar ese coste. Aplicado a los Impuestos Especiales españoles, el argumento es directo: quien conduce bajo los efectos del alcohol impone un riesgo a terceros que su propio precio de compra no refleja; quien fuma genera un coste sanitario público, y un coste de tabaquismo pasivo sobre quienes le rodean, que tampoco paga en el precio de la cajetilla; quien consume hidrocarburos contribuye a la contaminación atmosférica y al cambio climático, costes que recaen sobre el conjunto de la sociedad y no solo sobre quien repostó el depósito.

### 1.3. La crítica que el propio documento de referencia dejó apuntada sin desarrollar: la inelasticidad de la demanda

[ICEX-CECO] El documento original señala, en una frase suelta que nunca llega a desarrollar, que "algunos autores ponen en duda esta correlación" pigouviana, "por cuanto la demanda de estos productos tiende a ser inelástica". [Seguro] La objeción, formulada con precisión, es la siguiente: si la demanda de tabaco, alcohol o gasolina apenas responde al precio —porque se trata, en gran medida, de bienes de consumo habitual o adictivo, poco sustituibles—, entonces el impuesto especial cumple su función recaudatoria con gran eficacia, pero fracasa en su función correctora: el consumo apenas se reduce, y el impuesto se convierte, de facto, en un tributo puramente recaudatorio disfrazado de corrección de externalidad, cargando además —como ya vimos con el ejemplo mexicano de las bebidas azucaradas en el Tema 13— de forma desproporcionada sobre los hogares de renta baja, que dedican una proporción mayor de su presupuesto a estos bienes.

### 1.4. La respuesta moderna: la internalidad, y por qué la inelasticidad no invalida el argumento sino que lo refuerza

[Seguro] Esta objeción, formulada así, es exactamente la que la literatura de economía del comportamiento de las dos últimas décadas ha sometido a una revisión que conviene traer aquí con todo detalle, porque cambia por completo la conclusión. Jonathan Gruber y Botond Köszegi ("Is Addiction 'Rational'? Theory and Evidence", Quarterly Journal of Economics, 2001) parten de una evidencia empírica sólida: los fumadores son efectivamente previsores en sus decisiones —reaccionan, de forma medible, a subidas de impuestos especiales que ya han sido aprobadas legislativamente pero que todavía no han entrado en vigor, ajustando su consumo por anticipado—, lo que descarta la idea de que actúen de forma puramente impulsiva. Pero, simultáneamente, incorporan la evidencia igualmente sólida de que las preferencias sobre el tabaco son temporalmente inconsistentes: el fumador de hoy no valora sus decisiones futuras de la misma manera racional y estable que predeciría un modelo de descuento exponencial estándar, y con frecuencia ni siquiera reconoce con precisión lo difícil que le resultará dejarlo, buscando activamente mecanismos externos de autocontrol que le ayuden a conseguirlo. [Seguro] La implicación normativa de este modelo es radicalmente distinta de la del argumento pigouviano puro: el impuesto óptimo sobre el tabaco no debe calibrarse solo en función de la externalidad que el fumador impone sobre terceros, sino también en función de la internalidad —el daño que el propio fumador se impone a sí mismo, a su "yo futuro", precisamente por ese fallo de autocontrol—. Gruber y Köszegi estiman que esta internalidad asciende a en torno a 30 dólares por cajetilla, una cifra cien veces superior a la externalidad estimada sobre terceros, y concluyen que el impuesto óptimo debería ser, bajo su formulación, al menos un dólar por cajetilla más alto de lo que predeciría el modelo de adicción puramente racional. [Probable] Esto invierte por completo el sentido de la crítica de la inelasticidad que el documento de referencia dejaba sin resolver: si la demanda apenas responde al precio porque el consumidor está atrapado en un problema de autocontrol que él mismo reconocería como tal si se le preguntara directamente, entonces la inelasticidad no es un argumento contra el impuesto, sino la prueba misma de que el problema que el impuesto pretende corregir —la incapacidad de dejar de fumar aun deseándolo— es real y de una magnitud considerable.

[Probable] Conviene, como en cada debate de este tema, no cerrar la cuestión sin su contrapunto: Douglas Bernheim y Antonio Rangel ("Addiction and Cue-Triggered Decision Processes", American Economic Review, 2004) ofrecen un modelo alternativo en el que el consumo adictivo se desencadena por estímulos contextuales ("cue-triggered") más que por una genuina falta de autocontrol deliberativo, y concluyen que, bajo su formulación, los impuestos sobre bienes adictivos pueden no generar ningún cambio de comportamiento real y, en consecuencia, resultar reductores de bienestar antes que correctores de él —porque penalizan también al consumo "racional" y contextualmente no problemático, sin conseguir alterar el patrón de consumo verdaderamente compulsivo—. Ted O'Donoghue y Matthew Rabin ("Optimal Sin Taxes", Journal of Public Economics, 2006) intentan tender un puente entre ambas posiciones, formalizando cómo debería calibrarse un impuesto óptimo cuando la población se compone de una mezcla de consumidores plenamente racionales —para quienes el impuesto es una distorsión pura— y consumidores con problemas de autocontrol —para quienes el impuesto es, en cierto sentido, bienvenido—, concluyendo que el diseño óptimo depende de forma crítica de la proporción exacta de cada tipo dentro de la población, un parámetro que la evidencia empírica disponible no permite fijar con total precisión.

### 1.5. Marco jurídico y clasificación de las figuras españolas

[ICEX-CECO] Sentado el fundamento teórico, la regulación básica española es la Ley 38/1992, de 28 de diciembre, de Impuestos Especiales, y su Reglamento, aprobado por Real Decreto 1165/1995. Los Impuestos Especiales se clasifican en tres grandes grupos: los impuestos especiales de fabricación —que comprenden, a su vez, los impuestos sobre el alcohol y las bebidas alcohólicas (cuatro figuras distintas: alcohol y bebidas derivadas, cerveza, vino y bebidas fermentadas, y productos intermedios), el impuesto sobre hidrocarburos, el impuesto sobre las labores del tabaco, y el impuesto sobre la electricidad—; el impuesto especial sobre determinados medios de transporte; y el impuesto sobre el carbón. Se exigen en todo el territorio español excepto Canarias, Ceuta y Melilla, con las excepciones puntuales ya señaladas (cerveza, productos intermedios y alcohol/bebidas derivadas sí se exigen en Canarias; labores del tabaco sí se exige en Ceuta y Melilla; electricidad se exige en todo el territorio sin excepción). La recaudación de los impuestos especiales de fabricación está cedida a las Comunidades Autónomas en un 58%, salvo el de electricidad, cedido al 100%; el impuesto sobre medios de transporte también se cede al 100%. Conviene señalar, porque conecta con el argumento de armonización que ha atravesado todo este tema, que electricidad y medios de transporte son las dos únicas figuras de este bloque sin legislación armonizada a nivel de la Unión Europea — a diferencia del resto, que sí responden a Directivas comunitarias específicas, herederas de la misma lógica de armonización del mercado interior que dio origen al propio IVA.

### 1.6. Características comunes: por qué son estructuralmente distintos del IVA

[ICEX-CECO] Los Impuestos Especiales comparten cuatro rasgos que conviene contrastar, uno por uno, con las características del IVA que fijamos en el Bloque 1, porque la comparación es la mejor forma de fijarlos con precisión. Son indirectos y objetivos, como el IVA. Son instantáneos, como el IVA. Pero son, a diferencia del IVA, monofásicos: gravan una única fase —la fabricación o la importación—, no toda la cadena de producción y distribución. [Probable] Esta diferencia no es casual: recuerda que, en el Bloque 1 del Tema 13, ya establecimos que el diseño monofásico es preferible, precisamente, cuando se pretende concentrar el control en un número reducido de operadores fácilmente identificables —los fabricantes de cerveza, tabaco o hidrocarburos en España se cuentan por decenas o, como mucho, cientos, no por los cientos de miles de puntos de venta minorista que justificarían un diseño plurifásico—. El propio hecho imponible de fabricación o importación, con sus regímenes suspensivos que evitan la sujeción prematura mientras el producto permanece bajo control fiscal —en fábrica, en depósito fiscal, o en tránsito entre Estados miembros—, y la repercusión obligatoria que traslada la carga hasta el consumidor final, garantizan que, pese a gravarse en una única fase, la carga económica termine recayendo, igual que en el IVA, sobre quien finalmente consume el producto.


## 2. Los impuestos sobre el alcohol y las bebidas alcohólicas

### 2.1. Cuatro figuras, un mismo marco común

[ICEX-CECO] Bajo esta rúbrica conviven, en realidad, cuatro impuestos distintos, cada uno con su propio hecho imponible, su propia base imponible y su propia tarifa, aunque compartan reglas comunes de gestión: el impuesto sobre el alcohol y bebidas derivadas (alcohol etílico y bebidas espirituosas), el impuesto sobre la cerveza, el impuesto sobre el vino y las bebidas fermentadas, y el impuesto sobre los productos intermedios (bebidas con graduación entre el 1,2% y el 22% vol. no encuadrables en las categorías anteriores, como el jerez, el vermú o el oporto). Comparten, además, un conjunto de devoluciones específicas: por la utilización del alcohol en aromatizantes para productos alimenticios y bebidas analcohólicas; por su incorporación a productos alimenticios rellenos, con límites de contenido alcohólico según el tipo de producto; y por la devolución o destrucción, bajo control administrativo, de bebidas alcohólicas ya no aptas para el consumo humano.

### 2.2. El impuesto sobre el alcohol y bebidas derivadas

[ICEX-CECO] Grava la fabricación, importación e introducción en el territorio interno de alcohol etílico o etanol, total o parcialmente desnaturalizado, y de las bebidas derivadas aptas para consumo humano que contienen ese alcohol —no se exige en Ceuta y Melilla—. Además de las exenciones generales ya comunes a todos los impuestos especiales, se exime específicamente el alcohol no destinado a consumo humano o desnaturalizado, y el destinado a fabricación de medicamentos, investigación científica o fines docentes. La base imponible es el volumen de alcohol puro contenido, medido a 20ºC, y el tipo se expresa como un importe fijo en euros por hectolitro de alcohol puro —el criterio de medición más preciso posible desde el punto de vista pigouviano, porque grava directamente la magnitud —el etanol— que genera la externalidad, con independencia de en qué bebida concreta venga contenido—.

### 2.3. El impuesto sobre la cerveza

[ICEX-CECO] Grava la cerveza y las mezclas de cerveza con bebidas analcohólicas de graduación superior al 0,5% vol., tampoco exigible en Ceuta y Melilla. La base es el volumen en hectolitros de producto acabado a 20ºC, con una tarifa progresiva por tramos de graduación alcohólica: desde 0 €/hl para productos de graduación muy baja (hasta el 1,2% vol.) hasta importes crecientes conforme aumenta el contenido en alcohol. [Probable] Este diseño progresivo por graduación es, en sí mismo, un ejemplo bien construido de aplicación del principio pigouviano: cuanto mayor el contenido de alcohol, mayor el riesgo de externalidad por unidad de volumen consumido, y el tipo lo refleja de forma directa y proporcional.

### 2.4. El impuesto sobre el vino y las bebidas fermentadas: el gran contraste del sistema

[Seguro] Llegamos aquí al elemento más revelador de todo este apartado, y el que el documento de referencia se limita a constatar —"tipo impositivo de 0 €/hl para todos los productos desde su creación hasta hoy"— sin explicar por qué existe una asimetría tan llamativa dentro de un mismo bloque de impuestos, aparentemente construido sobre una misma lógica correctora. La razón no es económica: es jurídica y política, y hunde sus raíces en la propia arquitectura de la armonización europea de 1992 que ya trabajamos como marco general de todo el tema. La Directiva 92/84/CEE, que fija los tipos mínimos de accisa sobre el alcohol que todos los Estados miembros deben respetar como suelo, establece un mínimo positivo y sustancial para el alcohol etílico (550 ECU por hectolitro de alcohol puro, con salvedades históricas para algunos Estados) y para la cerveza, pero fija el tipo mínimo del vino en cero. [Probable] Esto significa que, a diferencia de lo que ocurre con el alcohol destilado o la cerveza, la propia normativa comunitaria no obliga a ningún Estado miembro a gravar el vino en absoluto — deja la decisión enteramente al arbitrio nacional, y España, como la inmensa mayoría de los grandes países productores de vino de la Unión —Francia, Italia, Portugal, además de Alemania para su producción vinícola—, ha optado sistemáticamente, desde la creación del impuesto en 1992, por mantener ese tipo en cero euros por hectolitro.

[Suposición] La explicación de esta asimetría no reside en ningún argumento pigouviano defendible —el alcohol contenido en una copa de vino genera, por unidad de etanol puro, exactamente la misma externalidad potencial que el alcohol contenido en una caña de cerveza o en una copa de licor—, sino en la capacidad de presión política del sector vitivinícola dentro de las negociaciones que condujeron a la Directiva de 1992, y que se ha mantenido intacta en las sucesivas décadas posteriores: los países del sur de Europa, con una industria vinícola de enorme peso económico, cultural y territorial —piensa en el propio mapa español de Denominaciones de Origen, de La Rioja a Ribera del Duero, de Jerez a Priorato—, consiguieron que el vino quedara excluido, de facto, del mismo tratamiento fiscal correctivo que sí se aplica al resto del alcohol. [Probable] Esta incoherencia no ha pasado desapercibida para la literatura de salud pública: la propia Organización Mundial de la Salud, en su Estrategia Mundial para Reducir el Uso Nocivo del Alcohol y en informes técnicos posteriores, ha recomendado reiteradamente que la fiscalidad sobre el alcohol se calibre en función del contenido puro de etanol, con independencia del tipo de bebida en que venga contenido —exactamente el criterio que sí aplican, dentro de España, tanto el impuesto sobre alcohol y bebidas derivadas como el de la cerveza, y que el impuesto sobre el vino incumple por completo al fijarse en cero—, señalando que la exención efectiva del vino frente a otras bebidas alcohólicas carece de justificación sanitaria y responde, en la práctica, a consideraciones de política agrícola y comercial ajenas al objetivo declarado de protección de la salud pública que en teoría legitima toda esta familia de impuestos.

### 2.5. El impuesto sobre productos intermedios y el carve-out de las Denominaciones de Origen

[ICEX-CECO] Grava los productos con graduación entre el 1,2% y el 22% vol. no encuadrables en cerveza ni en vino o bebidas fermentadas —el jerez, el vermú, el oporto, entre otros—, con base en el volumen a 20ºC y tipos escalonados según la graduación adquirida. [Seguro] Merece mención especial una excepción muy concreta: no está sujeta la fabricación de estos productos cuando ostenten las denominaciones de origen Moriles-Montilla, Tarragona, Priorato y Terra Alta, siempre que la adición de alcohol no incremente su graduación en más de un punto porcentual y se cumplan las condiciones reglamentarias correspondientes. [Probable] Esta excepción, aparentemente anecdótica por el reducido número de denominaciones afectadas, es en realidad la manifestación más nítida de la misma lógica de protección de la industria vitivinícola tradicional que ya identificamos con el tipo cero del vino: estas denominaciones producen, mediante métodos de elaboración centenarios, vinos generosos o de licor —el sistema de criaderas y soleras del Montilla-Moriles, por ejemplo— que técnicamente requieren una fortificación alcohólica mínima para alcanzar su perfil característico, y el legislador ha preferido eximir esa fortificación tradicional y acotada antes que someter a estas producciones artesanales, protegidas por su propia denominación geográfica, al mismo régimen que un producto intermedio industrial de gran escala.

### 2.6. Instrumentos alternativos que España no ha adoptado: el precio mínimo por unidad de alcohol

[Probable] Cierro este apartado señalando, como referencia comparada de actualidad, que el impuesto especial no es el único instrumento disponible para corregir la externalidad del consumo de alcohol, y que otros países han optado por herramientas distintas que España no ha incorporado a su arsenal de política fiscal. Escocia, desde 2018, aplica un precio mínimo por unidad de alcohol (Minimum Unit Pricing), que fija un suelo de precio por gramo de alcohol puro con independencia del tipo de bebida o de su coste de producción — un instrumento que, a diferencia del impuesto especial, no genera ingreso fiscal alguno para el Estado (la diferencia entre el precio de coste y el precio mínimo la retiene el propio productor o distribuidor), pero que persigue exactamente el mismo objetivo correctivo con una ventaja adicional: impide que las bebidas alcohólicas más baratas y de mayor graduación —históricamente, las más asociadas al consumo problemático— se sigan vendiendo por debajo de un umbral mínimo, algo que un impuesto especial expresado en euros por hectolitro, aplicado de forma uniforme dentro de cada categoría, no garantiza con la misma precisión si el producto de partida ya es extremadamente barato. [Suposición] Que España no haya explorado este instrumento, pese a la evidencia empírica favorable que los estudios de evaluación del caso escocés han ido documentando en términos de reducción de la mortalidad asociada al alcohol, es coherente con el patrón que ya identificamos al hablar del vino: en un país con una industria de bebidas alcohólicas de gran peso económico y territorial, cualquier instrumento de política fiscal que amenace con encarecer de forma sustancial ciertos segmentos de ese mercado se enfrenta a una resistencia política estructural que trasciende, con mucho, el debate puramente técnico sobre cuál es el instrumento correctivo óptimo.


## 3. El impuesto sobre hidrocarburos

### 3.1. Marco básico

**[ICEX-CECO]** Grava la fabricación e importación, en el ámbito territorial interno, de gasolinas con y sin plomo, gasóleos, fuelóleos, gases licuados del petróleo, metano, queroseno y una amplia gama de productos que adquieren la calificación fiscal de hidrocarburos cuando se destinan a usarse como combustible o carburante. **Se aplica únicamente en la Península y las Islas Baleares**, quedando excluidos Canarias, Ceuta y Melilla. La base imponible se mide en volumen (miles de litros a 15ºC) o, para determinados productos, en peso (toneladas métricas) o en poder energético (gigajulios). Los tipos se organizan en dos tarifas —la primera, para los productos de uso más habitual como combustible o carburante; la segunda, para el resto—, con tipos reducidos específicos para gasóleo de calefacción, gasóleo agrícola, gas natural para usos no combustibles y queroseno para usos distintos del transporte.

[Seguro] No repito aquí el mecanismo del fraude de operador desaparecido aplicado a este sector —ya lo desarrollamos con todo el detalle que merece en el Bloque 1, con los casos reales de Copérnico y Pamplinas Stars—; recuerda, simplemente, que el hidrocarburo es, junto al tabaco, la figura de mayor recaudación de todo este bloque, y por tanto la que concentra el mayor atractivo para el fraude organizado a gran escala.

### 3.2. El episodio del "céntimo sanitario": cuando una Comunidad Autónoma inventa un impuesto que Europa termina anulando

[Seguro] El episodio más dramático y mejor documentado de toda la historia reciente de este impuesto no afecta al propio Impuesto sobre Hidrocarburos estatal, sino a un tributo que se le superpuso durante una década: el **Impuesto sobre las Ventas Minoristas de Determinados Hidrocarburos (IVMDH)**, conocido popularmente como **"céntimo sanitario"**, creado por el art. 9 de la Ley 24/2001 y modificado por el art. 7 de la Ley 53/2002, vigente entre 2002 y 2012. Se trataba de un tributo autonómico, de naturaleza indirecta, que gravaba en fase única las ventas minoristas de determinados hidrocarburos, cuya recaudación las Comunidades Autónomas tenían la facultad de destinar, formalmente, a la financiación de gastos sanitarios —de ahí su apodo—.

[Seguro] Una empresa de transporte de mercancías catalana, **Transportes Jordi Besora, S.L.**, que había abonado 45.632,38 euros en concepto de este impuesto entre 2005 y 2008, lo consideró contrario a la Directiva europea sobre impuestos especiales y reclamó su devolución, dando origen a una cuestión prejudicial elevada por el Tribunal Superior de Justicia de Cataluña ante el TJUE. [Seguro] La **sentencia del TJUE de 27 de febrero de 2014** (asunto C-82/12) declaró el IVMDH **contrario al Derecho de la Unión**, con un razonamiento de gran interés técnico: la Directiva sobre impuestos especiales permite a los Estados miembros mantener impuestos indirectos adicionales sobre productos ya sujetos a un impuesto especial armonizado, pero **solo si persiguen una finalidad específica**, distinta de la meramente presupuestaria. El Tribunal consideró que el "céntimo sanitario", pese a su etiqueta y a su vinculación formal a la financiación sanitaria, **no cumplía ese requisito**: la mera afectación contable de lo recaudado a un fin determinado no basta para constituir una "finalidad específica" en el sentido exigido por la Directiva, si la estructura misma del impuesto —su hecho imponible, su base, su tipo— no está diseñada para desincentivar, de forma directa y verificable, el consumo del producto gravado por esa razón sanitaria concreta.

[Seguro] Lo más severo de la sentencia, y lo que conviene retener como el dato de mayor calado jurídico de todo el episodio, es que el Tribunal **rechazó expresamente limitar los efectos temporales de su fallo** —algo que sí suele conceder en casos de gran impacto fiscal, precisamente para proteger la seguridad jurídica y las finanzas públicas del Estado afectado—, razonando que el Gobierno español y la Generalitat de Catalunya **no habían actuado de buena fe** al mantener el impuesto en vigor durante más de diez años pese a las dudas, ya conocidas, sobre su compatibilidad con el Derecho europeo. [Probable] La consecuencia práctica de esta negativa a limitar los efectos temporales fue una **oleada masiva de reclamaciones de devolución** por parte de miles de contribuyentes —principalmente estaciones de servicio, que habían actuado como sujetos pasivos formales del impuesto, pero también empresas transportistas y otros consumidores finales que habían soportado la repercusión—, que se prolongó durante años a través de dos vías procesales distintas y en competencia —la rectificación de autoliquidaciones y devolución de ingresos indebidos, de un lado, y la exigencia de responsabilidad patrimonial del Estado legislador, de otro—, generando una litigiosidad de gran volumen y una factura fiscal muy sustancial para las arcas autonómicas y estatales por un tributo que, durante toda su vigencia, se había presentado ante la ciudadanía como una medida de protección de la sanidad pública.

[Suposición] El episodio del céntimo sanitario merece leerse, dentro de todo lo que hemos construido en este tema, como una advertencia de alcance general: incluso dentro de un sistema de imposición especial tan intensamente armonizado como el europeo, un Estado miembro —o, como en este caso, una Comunidad Autónoma con capacidad normativa cedida— puede diseñar un tributo que, formalmente, pretenda encajar en las excepciones que la propia armonización permite, y sin embargo fracasar en ese intento si la finalidad declarada no se traduce en un diseño técnico coherente con ella. Es, en cierto sentido, el mismo tipo de tensión entre la letra de la norma y su verdadera función económica que ya trabajamos, con otro objeto, a propósito de la interpretación estricta de las exenciones de IVA en el Bloque 2.

### 3.3. Un debate de actualidad no resuelto: la brecha fiscal entre gasóleo y gasolina

[Probable] Cierro con un debate que sigue abierto y que conecta directamente con la propia agenda de fiscalidad medioambiental del Libro Blanco de 2022 que ya trabajamos en el Bloque 2. Durante décadas, la tarifa española —como la de la mayoría de los países europeos— ha gravado el **gasóleo** a un tipo sensiblemente **inferior** al de la **gasolina**, una diferencia que en su origen respondía a razones de política industrial y de apoyo al transporte profesional y agrícola —el gasóleo se asociaba tradicionalmente al transporte de mercancías y a la maquinaria agrícola, sectores que el legislador quería proteger frente a un mayor coste energético—. [Probable] El escándalo internacional conocido como *"Dieselgate"* —la manipulación por parte de Volkswagen, revelada en 2015, de los sistemas de control de emisiones de sus motores diésel para superar fraudulentamente las pruebas de homologación— puso en cuestión, a escala europea, la propia premisa medioambiental que durante años había acompañado a esa ventaja fiscal: el diésel, lejos de ser la alternativa más limpia que su menor tributación parecía sugerir, resultó emitir, en condiciones reales de conducción, niveles de óxidos de nitrógeno muy superiores a los declarados, con un impacto documentado sobre la calidad del aire urbano y la salud respiratoria comparable o superior al de la gasolina. El propio Comité de Personas Expertas del Libro Blanco de 2022, dentro de su capítulo de fiscalidad medioambiental —con una estimación de recaudación adicional potencial de entre 5.941 y 15.023 millones de euros para el conjunto de las propuestas verdes—, señaló el transporte como uno de sus dos ámbitos de actuación prioritaria, apuntando implícitamente a la necesidad de revisar esta diferencia histórica de trato fiscal entre ambos combustibles. [Suposición] Que esta convergencia no se haya producido todavía de forma sustancial en España, pese al consenso técnico cada vez más amplio sobre su justificación medioambiental, es coherente con el mismo patrón de secuenciación política que ya identificamos al cerrar el Bloque 2: el diagnóstico técnico está disponible desde hace años, pero su traducción en una subida efectiva de un impuesto que afecta de forma directa al coste del transporte de mercancías y al presupuesto de millones de conductores particulares sigue topando con la misma resistencia política que cualquier reforma fiscal de calado en España.


\subsection{Impuesto sobre las labores del tabaco}
Grava la fabricación e importación, en territorio peninsular español e Islas Baleares, de cigarros y cigarritos, cigarrillos, picadura para liar y demás tabacos de fumar. La base imponible se determina de dos formas según el tipo aplicable: para los **tipos proporcionales**, el valor de las labores según su precio máximo de venta al público en expendedurías, impuestos incluidos; para los **tipos específicos**, el número de unidades. Los tipos pueden ser **proporcionales** (porcentaje sobre el PVP), **específicos** (importe fijo por unidad o peso), o ambos **simultáneamente**, con un **tipo mínimo** adicional que actúa como suelo de tributación cuando la combinación de ambos no lo alcanza.

### 4.2. Por qué el tabaco combina los tres mecanismos: Suits y Musgrave aplicados al caso real

[Probable] Este diseño híbrido —proporcional, específico y mínimo, los tres a la vez— no es un capricho de complejidad normativa: es la aplicación práctica, casi de manual, del resultado de **Suits y Musgrave** que ya trabajamos en el Bloque 1 de este mismo tema. Recuerda su conclusión: bajo competencia imperfecta —y el mercado del tabaco, dominado por un número reducido de grandes fabricantes multinacionales, es precisamente ese caso—, un impuesto puramente **específico** genera una subida de precio y una reducción de cantidad mayores que uno **ad valorem** de recaudación equivalente, porque el ad valorem reduce el margen del propio fabricante y modera, por esa vía, su propio incentivo a subir el precio. [Suposición] Pero un tipo puramente ad valorem tiene, a su vez, un defecto que el legislador tributario evidentemente quiso evitar: al ser proporcional al precio, favorece de forma desproporcionada a las marcas más baratas —el fabricante de una marca de gama baja paga, en términos absolutos, mucho menos impuesto por cada cigarrillo que el de una marca premium, aunque el daño sanitario por cigarrillo consumido sea idéntico—, lo que puede empujar el consumo hacia esas marcas de menor precio en lugar de reducirlo en términos agregados. El **tipo específico**, al gravar por unidad con independencia del precio, corrige exactamente ese sesgo, igualando la carga fiscal absoluta entre marcas caras y baratas. Combinar ambos —y blindar el conjunto con un **tipo mínimo** que garantiza que ninguna combinación de precio y volumen pueda escapar a un suelo de tributación— es, por tanto, la respuesta de diseño que sintetiza las ventajas de cada uno de los dos instrumentos que Suits y Musgrave analizaron por separado, neutralizando el defecto específico de cada uno con la presencia del otro.

### 4.3. El comercio ilícito: la confirmación empírica más nítida de Adam Smith

[Seguro] Ya trabajamos, en el Tema 13, la crítica clásica de Adam Smith sobre cómo un impuesto especial excesivamente alto sobre un bien fácilmente transportable genera, casi inevitablemente, un incentivo al contrabando. El tabaco es, con diferencia, el mejor caso de estudio contemporáneo de esa misma dinámica, porque España dispone de datos empíricos abundantes y bien documentados sobre su magnitud exacta a lo largo del tiempo. [Probable] Según el informe AFI (Analistas Financieros Internacionales) de mayo de 2013, el comercio ilícito de tabaco en España se **triplicó** en pocos años, pasando de representar el 2,2% del mercado a el 7,5% en 2012, con una pérdida fiscal estimada, en el escenario más desfavorable, de hasta **1.631 millones de euros** anuales. Los sucesivos informes *Project SUN* de KPMG documentan después una senda de mejora sostenida: el consumo ilícito bajó al 4,6% en 2016 —Rocío Ingelmo, directora de Asuntos Corporativos de Altadis, defendió en su propia tesis doctoral la existencia de una relación directa y cuantificable entre las subidas de impuestos, el consiguiente diferencial de precio con los mercados vecinos, y el aumento de las tasas de contrabando, un efecto que además se **acentúa en periodos de mayor desempleo**, cuando la elasticidad precio-demanda de los fumadores aumenta de forma significativa—, y continuó bajando hasta el **3,3% en 2021**, con una pérdida fiscal ya reducida a 248 millones de euros, el nivel más bajo desde 2010.

[Probable] Conviene, con la misma honestidad metodológica que hemos aplicado a cada fuente de este tema, señalar una precisión sobre el origen de estos datos: los informes *Project SUN* y su sucesor *Project Stella* están **encargados y financiados por Philip Morris International**, la mayor tabacalera multinacional del mundo, en el marco de un acuerdo comercial que puso fin a acciones legales que la propia Unión Europea había emprendido contra la compañía por su implicación histórica en el comercio ilícito de tabaco. Organizaciones independientes de control del tabaquismo han señalado que estos informes tienden a etiquetar como "otro contrabando y falsificación" una parte sustancial del tabaco ilícito que, según datos de incautaciones globales, procedería en realidad de la propia industria tabacalera —hasta dos tercios del mercado ilícito global, según esas fuentes—, lo que sugiere una posible infraestimación sistemática del papel de los propios fabricantes en la dinámica del comercio ilícito que sus informes, financiados por ellos mismos, describen.

### 4.4. Una singularidad geográfica española: tres precios, un solo país

[Seguro] Un rasgo distintivo del caso español, poco replicado en el resto de la Unión Europea, es que **España es el único país de la UE con tres niveles de precio del tabaco** dentro de su propia jurisdicción: uno para la Península y Baleares, otro fijado específicamente para Ceuta y Melilla, y un tercero, libre de impuesto especial, para Canarias. [Probable] Esta triple fragmentación interna genera, por sí sola, incentivos de arbitraje fiscal **dentro del propio territorio nacional**, sin necesidad siquiera de cruzar una frontera internacional —el mismo tipo de incentivo que ya identificamos, en otro contexto, al hablar de las tramas de fraude de hidrocarburos—. A esto se suma la proximidad de **Gibraltar**, con precios sensiblemente inferiores a los peninsulares, convertido en uno de los focos de contrabando de tabaco más señalados de toda Europa según diversos estudios sobre comercio ilícito en la Unión. [Suposición] Y, de forma menos conocida pero igualmente reveladora, España no es solo receptora neta de tabaco ilícito: según el propio informe Project SUN, España es también **exportadora neta de tabaco con impuesto ya pagado** hacia Francia y el Reino Unido, cuyos precios domésticos prácticamente duplican los españoles —un flujo de arbitraje fiscal en sentido inverso al que domina el relato público habitual sobre este fenómeno, y que conviene conocer para no simplificar en exceso la dirección del problema—.

### 4.5. La frontera regulatoria más reciente: el vapeo y el tabaco calentado

[Suposición] Cierro con un debate de plena actualidad que la Ley 38/1992, redactada mucho antes de que estos productos existieran comercialmente, no anticipó: los **cigarrillos electrónicos y líquidos de vapeo**, y en menor medida el **tabaco calentado**, ocupan hoy un vacío regulatorio parcial dentro del sistema de impuestos especiales, al no encajar con naturalidad en la definición tradicional de "labores del tabaco" —los líquidos de vapeo, en particular, no contienen hoja de tabaco alguna—. Varios países de la Unión Europea han avanzado ya hacia la creación de gravámenes específicos para estos productos, precisamente para evitar que se conviertan en una vía de escape fiscal frente al tabaco convencional cada vez más gravado, mientras que en España el debate sobre cómo encajarlos dentro del marco de los Impuestos Especiales —o si crear, en su lugar, una figura tributaria enteramente nueva— sigue sin resolverse de forma definitiva, precisamente el tipo de laguna regulatoria que conviene que sepas identificar como reflejo de una norma diseñada para un mercado del tabaco que la propia innovación del sector ha terminado por desbordar.


## 5. El impuesto sobre la electricidad y el impuesto sobre el carbón

### 5.1. El impuesto sobre la electricidad

**[ICEX-CECO]** Impuesto especial de naturaleza indirecta que recae sobre el consumo de electricidad, gravando en fase única el suministro para consumo del adquirente y el autoconsumo por el propio productor —con exclusión, por razones de proporcionalidad administrativa que ya conoces de otros contextos, del autoconsumo de instalaciones de potencia no superior a 100 kW—. La base imponible es la misma que se habría determinado a efectos de IVA, excluida la cuota del propio impuesto especial; el tipo es del **5,11269632%**, con una garantía de imposición mínima —0,5 €/MWh para usos industriales, embarcaciones no de recreo y transporte ferroviario; 1 €/MWh para el resto—.

[Seguro] El origen de estas dos cifras mínimas no es una elección discrecional del legislador español: proceden directamente de la **Directiva 2003/96/CE del Consejo, de 27 de octubre de 2003**, por la que se reestructura el régimen comunitario de imposición de los productos energéticos y la electricidad, que fija esos mismos umbrales —0,5 y 1 euro por megavatio-hora— como suelo mínimo obligatorio para todos los Estados miembros. Es, otra vez, la misma lógica de armonización con suelos mínimos que ya vimos operar, con cifras distintas, en la Directiva sobre el alcohol.

[Seguro] El tipo del 5,11269632% —una cifra con siete decimales que rara vez se explica, y que conviene retener como ejemplo de hasta qué punto la técnica tributaria puede parecer arbitraria sin serlo— fue, además, **reducido de forma temporal al 0,5%** desde 2021 hasta el 31 de diciembre de 2022, mediante una cadena de Reales Decretos-leyes sucesivos (12/2021, 17/2021, 29/2021, 6/2022, 11/2022) que fueron prorrogando la medida trimestre a trimestre —exactamente el mismo patrón de "medida temporal que se prorroga indefinidamente" que ya documentamos con el IVA de los alimentos básicos en el Bloque 2—, como respuesta a la crisis de precios energéticos desatada por la invasión rusa de Ucrania.

[Seguro] Conviene, además, situar este impuesto dentro de un panorama fiscal energético español más amplio del que el propio documento de referencia no da cuenta: junto al Impuesto Especial sobre la Electricidad que grava el **consumo**, existe un tributo distinto, el **Impuesto sobre el Valor de la Producción de la Energía Eléctrica (IVPEE)**, creado por la **Ley 15/2012**, que grava, en cambio, a los **productores** de electricidad, justificado en su exposición de motivos por las inversiones en redes de transporte y distribución que su actividad genera, por sus efectos medioambientales, y por los costes de garantía de suministro asociados. [Seguro] Este segundo impuesto fue sometido a un intenso escrutinio de compatibilidad con el Derecho de la Unión, en un episodio que conviene contrastar con el fracaso del céntimo sanitario que ya narramos a propósito de los hidrocarburos: el **TJUE, en su sentencia de 3 de marzo de 2021** (asunto C-220/19, *Oliva Park*), a diferencia de lo ocurrido con el IVMDH, **declaró el IVPEE compatible** con el Derecho de la Unión, descartando las dudas planteadas sobre su encaje con la Directiva armonizadora de impuestos especiales, con la Directiva de fomento de energías renovables, y con el principio de libre competencia en el mercado interior de la electricidad. [Probable] Dos tributos energéticos españoles, litigados ante el mismo tribunal por razones de compatibilidad europea, con desenlaces opuestos: uno anulado por carecer de finalidad específica genuina, el otro validado por superar ese mismo escrutinio — una buena ilustración, para cerrar este apartado, de que el fracaso del céntimo sanitario no revela una hostilidad genérica del Derecho europeo hacia la fiscalidad energética española, sino una exigencia de coherencia técnica entre la finalidad declarada de cada tributo y su diseño material real, exigencia que el IVPEE, a diferencia del IVMDH, sí superó.

### 5.2. El impuesto sobre el carbón

**[ICEX-CECO]** Grava la **puesta a consumo** de carbón —la primera venta o entrega tras su producción, extracción, importación o adquisición intracomunitaria, así como su autoconsumo—, con base en su poder energético expresado en gigajulios, y dos tipos: **0,15 €/GJ** para usos profesionales distintos de la generación eléctrica, y **0,65 €/GJ** para el resto de usos.

[Probable] Lo más revelador de este impuesto, precisamente por su carácter marginal dentro de la recaudación total del bloque, es lo que su propio diseño **excluye**: el carbón destinado a la **generación y cogeneración eléctrica** —con diferencia, su uso más relevante y más intensivo en emisiones— queda fuera del ámbito de aplicación del tipo general, precisamente porque esa combustión ya está sometida a un instrumento de precio del carbono completamente distinto y de mayor alcance: el **Régimen de Comercio de Derechos de Emisión de la Unión Europea (EU-ETS)**, el mismo mecanismo que ya trabajamos con todo detalle en el Bloque 1 de este tema, a propósito del fraude carrusel de 2008-2009. [Suposición] Esta exclusión no es casual: gravar dos veces la misma tonelada de carbón quemada en una central eléctrica —una vez mediante el impuesto especial nacional, y otra mediante el precio del derecho de emisión europeo— constituiría una duplicación de instrumentos correctivos sobre la misma externalidad, precisamente el tipo de incoherencia de diseño que hemos aprendido a identificar a lo largo de todo este bloque. El impuesto especial sobre el carbón queda, así, reservado a un papel **residual y complementario**: alcanza únicamente los usos industriales de combustión directa —hornos, procesos de reducción química— que quedan fuera del ámbito del EU-ETS, explicando por qué su recaudación es, con gran diferencia, la más reducida de todo el catálogo de impuestos especiales de fabricación que hemos recorrido en este bloque. [Probable] A esto se suma un dato de contexto que conviene retener como cierre: el cierre casi completo del parque de centrales térmicas de carbón español, culminado hacia 2020 bajo la presión combinada de la normativa medioambiental europea y de un precio del derecho de emisión del EU-ETS que hizo, durante años, económicamente inviable la generación eléctrica a partir de carbón frente a otras fuentes, ha convertido a este impuesto en una figura cada vez más vestigial dentro del sistema español de imposición especial —una figura que sigue formalmente vigente, pero cuya base imponible real se ha ido reduciendo, año tras año, al ritmo exacto de la propia transición energética que otros instrumentos, ajenos a este impuesto concreto, han terminado por imponer.



## 6. El impuesto especial sobre determinados medios de transporte

### 6.1. Marco básico y la cita de la exposición de motivos

**[ICEX-CECO]** Grava la primera matriculación definitiva en España de vehículos automóviles nuevos o usados, de embarcaciones y buques de recreo de más de 7,5 metros de eslora, y de aeronaves provistas de motor mecánico, con un catálogo de exenciones que ya recorrimos en el Bloque C del documento de referencia —taxis, autoescuelas, vehículos de alquiler, personas con discapacidad, uso diplomático, embarcaciones a remo o vela olímpica, aeronaves de titularidad pública o dedicadas a la enseñanza—. La propia exposición de motivos de la norma, que ya citamos en su momento, formula con precisión el fundamento del tributo: no responde "exclusivamente al gravamen de la capacidad contributiva puesta de manifiesto en su adquisición", sino a "la consideración adicional de las implicaciones de su uso en la producción de costes sociales específicos en el ámbito de la sanidad, las infraestructuras o el medio ambiente" — el mismo argumento pigouviano que ha vertebrado todo este bloque, aplicado ahora al momento concreto de la matriculación en lugar de al consumo periódico del combustible.

### 6.2. La reforma de 2007: de la cilindrada a las emisiones de CO2

[Seguro] El episodio más relevante en la historia reciente de este impuesto es su reestructuración completa mediante la disposición final octava de la **Ley 34/2007, de 15 de noviembre, de calidad del aire y protección de la atmósfera**, en vigor desde el **1 de enero de 2008**. Hasta entonces, la base de tributación había sido la **cilindrada** del motor —un criterio que, entre otras cosas, favorecía de forma indirecta y no deliberada a los motores diésel frente a los de gasolina, por su distinta relación entre cilindrada y potencia—. La reforma sustituyó ese criterio por uno directamente vinculado a las **emisiones de CO2 por kilómetro**, articulado en cuatro tramos: **exención total** para vehículos con emisiones inferiores a 120 g/km; **4,75%** sobre el valor del vehículo para el tramo de 121 a 160 g/km; **9,75%** para el tramo de 161 a 200 g/km; y **14,75%** para vehículos con emisiones iguales o superiores a 200 g/km —tramo en el que, además, se incorporaron por primera vez al impuesto las motos acuáticas y los quads, que hasta entonces no tributaban en absoluto—. [Seguro] El contexto que motivó la reforma es revelador: la Comisión Europea había fijado el objetivo de que la media de emisiones de los turismos matriculados en la Unión no superara los **120 g/km para 2012**, mientras que la media española de 2006 se situaba en **151,91 g/km**, y apenas el **15%** de los vehículos vendidos ese año habría cumplido el nuevo umbral exento —una brecha que explica por qué la reforma, lejos de ser un ajuste cosmético, tuvo un impacto recaudatorio y de comportamiento del mercado muy significativo desde su primer año de aplicación, según reflejan los propios datos de matriculaciones de la Agencia Tributaria del ejercicio 2008.

### 6.3. La ausencia de armonización europea, y el límite que impone el art. 110 del TFUE

**[ICEX-CECO]** Ya señalamos, al presentar el marco general de este bloque, que el impuesto sobre medios de transporte es, junto al de electricidad, una de las dos figuras **sin legislación armonizada** a nivel de la Unión Europea. [Probable] Esto no significa que el impuesto opere en un vacío jurídico europeo completo: aunque la Comisión Europea presentó en 2005 una propuesta de Directiva para armonizar la fiscalidad del automóvil en la Unión —que habría avanzado hacia la sustitución progresiva de los impuestos de matriculación por impuestos de circulación basados en emisiones—, la propuesta nunca alcanzó el consenso necesario en el Consejo y fue formalmente **retirada por la Comisión en 2012**, dejando la fiscalidad de la matriculación como una de las pocas áreas de imposición indirecta relevante en las que cada Estado miembro conserva plena libertad de diseño. [Seguro] Esa libertad, sin embargo, **no es absoluta**: el art. 110 del Tratado de Funcionamiento de la Unión Europea prohíbe a los Estados miembros gravar los productos de otros Estados miembros con tributos internos superiores a los que gravan productos nacionales similares, y el Tribunal de Justicia ha desarrollado, en una línea jurisprudencial muy consolidada —con sentencias recientes contra Portugal (C-169/20, 2 de septiembre de 2021) y contra Malta (C-694/22, 22 de febrero de 2024)—, una doctrina específica para los **vehículos usados importados**: cuando un impuesto de matriculación se incorpora al valor residual de un vehículo ya matriculado en el mercado nacional —porque ese valor residual refleja, de forma implícita, el impuesto ya pagado en su día—, el impuesto que grava un vehículo usado importado de otro Estado miembro debe tener en cuenta esa **depreciación** exacta, para no acabar gravando el vehículo importado por un importe superior al que ya lleva incorporado, de forma invisible, el vehículo usado nacional equivalente. [Suposición] Es un límite de coherencia técnica muy exigente, y que cualquier reforma futura del IEDMT español debería seguir respetando con precisión, precisamente el mismo tipo de vigilancia de neutralidad entre producto nacional e importado que ya vimos operar, con otro objeto y otro tribunal, a propósito de la interpretación estricta de las exenciones de IVA en el Bloque 2.

### 6.4. El futuro del impuesto: la erosión de base que la propia electrificación genera

[Suposición] Cierro este último apartado con el debate de mayor actualidad y el que, probablemente, determinará la evolución de esta figura en los próximos años. El diseño de 2007-2008, tan elogiado en su momento por vincular la carga fiscal directamente a las emisiones de CO2, contiene una fragilidad estructural que sus propios autores difícilmente podían anticipar hace casi dos décadas: un **vehículo eléctrico puro**, con emisiones de tubo de escape iguales a cero, queda automáticamente exento del impuesto bajo el primer tramo de la tarifa —una consecuencia perfectamente coherente con la lógica pigouviana del diseño, pero que, a medida que la cuota de mercado de los vehículos eléctricos crece de forma sostenida, va **erosionando progresivamente la base imponible** de un impuesto cuya arquitectura entera está construida sobre la premisa de que los vehículos que se matriculan emiten CO2 en cantidades variables pero nunca nulas. [Probable] Este mismo problema, replicado a escala mucho mayor en el impuesto sobre hidrocarburos —cuya recaudación depende, de forma todavía más directa, del consumo de combustibles fósiles que la propia transición energética pretende reducir hasta su desaparición—, ha generado en toda Europa un debate creciente sobre instrumentos sustitutivos de recaudación vinculados no ya al tipo de motor ni a las emisiones, sino directamente al **uso efectivo de la infraestructura viaria**: los llamados sistemas de **pago por kilómetro recorrido** (*road pricing* o *distance-based charging*), que sustituirían tanto la fiscalidad sobre el combustible como la de matriculación por un gravamen calculado en función de los kilómetros efectivamente circulados, con independencia de si el vehículo es eléctrico, híbrido o de combustión. [Probable] Los Países Bajos fueron pioneros en plantear esta transición ya en 2007, cuando el entonces ministro de Transporte, Camiel Eurlings, anunció un proyecto de gravamen por kilómetro recorrido —con registro por satélite— que debía sustituir gradualmente al impuesto de circulación, primero para camiones y después para el conjunto del parque automovilístico; el proyecto neerlandés terminaría abandonándose por la magnitud de la resistencia política y técnica que generó, pero la misma idea de fondo ha reaparecido, con distinto grado de desarrollo, en varios países europeos y en la propia agenda de reflexión de la Comisión Europea sobre el futuro de la fiscalidad del transporte en un escenario de electrificación masiva del parque —un debate que, a diferencia de los que hemos recorrido a lo largo de todo este bloque, no versa sobre corregir una externalidad ya identificada con un instrumento ya maduro, sino sobre cómo **reinventar la propia base imponible** de toda una familia de impuestos cuya premisa técnica original está siendo desmontada, año tras año, por el éxito mismo de la política medioambiental que esos impuestos pretendían impulsar.

## Valoración general del Bloque 3

### El hilo que conecta las seis figuras: el mismo compromiso, seis veces distinto

Recorridas las seis figuras, conviene detenerse a mirar el conjunto y no solo cada pieza por separado, porque un patrón se repite con una regularidad que merece nombrarse: en **cada una** de estas seis figuras, la justificación pigouviana declarada —corregir una externalidad— convive con una segunda fuerza, casi siempre más determinante en la práctica, que tuerce el diseño técnico hacia algo distinto de lo que la pura teoría recomendaría. En el alcohol, esa segunda fuerza es la **capacidad de presión política** de la industria vitivinícola, capaz de mantener el vino a tipo cero durante más de tres décadas pese a compartir, molécula por molécula, la misma externalidad que la cerveza o los espirituosos sí soportan. En los hidrocarburos, es la **tentación presupuestaria autonómica**, que llevó a Cataluña y al resto de Comunidades a inventar un "céntimo sanitario" cuya finalidad declarada no resistió el escrutinio europeo. En el tabaco, es el **comercio ilícito** que la propia intensidad recaudatoria del impuesto genera como efecto colateral, con datos que van de los 1.631 millones de pérdida estimada en el peor año a los 248 millones del más reciente. En la electricidad, es la **volatilidad de la política energética de emergencia**, que ha convertido un tipo nominal del 5,11269632% en una cifra casi teórica, sustituida durante años por reducciones "temporales" prorrogadas sin cesar. En el carbón, es la **obsolescencia sobrevenida**: un impuesto que sigue vigente sobre un combustible que la propia política europea de descarbonización, a través de un instrumento distinto —el EU-ETS—, ha conseguido casi expulsar del sistema eléctrico español. Y en los medios de transporte, es la **erosión de base por éxito de la propia política**: un diseño ejemplar en 2008 que la electrificación del parque automovilístico está vaciando de contenido recaudatorio precisamente porque está cumpliendo su función correctora mejor de lo previsto.

### Qué figura está mejor diseñada, y por qué no es la que más recauda

[Probable] Si tuvieras que jerarquizar estas seis figuras por coherencia interna entre su justificación teórica y su diseño técnico real, el **tabaco** sale, con diferencia, mejor parado: su estructura híbrida —proporcional, específico y mínimo combinados— es la aplicación más depurada de Suits y Musgrave de todo el bloque, y el argumento de la internalidad de Gruber y Köszegi le proporciona, además, un fundamento normativo que ni siquiera necesita apoyarse en la externalidad sobre terceros para justificar tipos elevados. El **alcohol**, en el extremo opuesto, es la figura con la incoherencia más flagrante: ningún manual de finanzas públicas explicaría satisfactoriamente por qué el etanol de una copa de vino escapa por completo a un impuesto que sí grava, con precisión creciente según la graduación, el etanol de una cerveza o de un licor. [Suposición] Y sin embargo, la recaudación no premia la coherencia: son los **hidrocarburos** —con su propia historia de litigios, fraude y debates de convergencia diésel-gasolina sin resolver— los que aportan, con gran diferencia, la mayor partida de todo el bloque, seguidos del tabaco; mientras que el carbón, probablemente la figura menos discutible desde un punto de vista puramente pigouviano dado el perfil de emisiones de este combustible, es hoy la de menor peso recaudatorio de las seis. La lección que se extrae de este contraste es la misma que ya anticipamos al hablar de la relación entre coherencia teórica y resultado práctico en el Bloque 2: **el diseño más elegante no es, casi nunca, el que más dinero recauda**, y cualquier reforma que persiga simultáneamente ambos objetivos —coherencia y recaudación— se enfrenta a una tensión que ningún comité de expertos, por bien fundamentado que esté, ha conseguido resolver del todo.

### La conexión final con Puviani: los impuestos especiales como el caso más puro de ilusión fiscal

[Probable] Cierro con una idea que conecta este bloque final con uno de los primeros conceptos que trabajamos en todo el tema, y que conviene subrayar como broche del recorrido completo: si el IVA, por su propio diseño de repercusión obligatoria y desglose explícito en cada factura, es relativamente **transparente** para quien lo paga —la hipótesis de Mill que verificamos empíricamente con Chetty, Looney y Kroft—, los Impuestos Especiales son, de las seis familias tributarias que hemos recorrido, la que mejor encarna la **teoría de la ilusión financiera de Puviani**: nadie que reposta gasolina, compra una cerveza o adquiere una cajetilla de tabaco ve, en el ticket de compra, un desglose que le diga cuánto de ese precio es impuesto especial y cuánto es coste del propio producto —a diferencia del IVA, que sí aparece desglosado con obligación legal—. El impuesto especial queda **enterrado dentro del precio final**, de una forma mucho más profunda que el propio IVA que hemos estudiado en los dos primeros bloques, y esa opacidad no es casual: es, precisamente, la que permite que estas seis figuras sostengan tipos efectivos extraordinariamente altos —recuerda que el tabaco llegaba a tipos efectivos superiores al 80-90% del precio final en los datos estadísticos del inicio de este tema— sin generar, en el debate público, una resistencia proporcional a su magnitud real.

### Cierre del Tema 20 completo

Con esta valoración queda cerrado el Bloque 3, y con él el tema en su totalidad. El recorrido ha trazado un arco deliberado: el Bloque 1 presentó el IVA en su forma más neutra y coherente, el diseño de Lauré funcionando exactamente como su teoría predice; el Bloque 2 introdujo las primeras grietas en esa neutralidad —exenciones y tipos diferenciados— y mostró que incluso esas grietas obedecen a una lógica identificable, aunque en tensión permanente con la eficiencia pura; y este Bloque 3 ha cerrado el recorrido con la forma más alejada de la neutralidad de todas las que ofrece la imposición indirecta española —la selectividad deliberada—, mostrando que ni siquiera aquí, en el terreno donde el legislador declara con mayor claridad sus intenciones correctoras, el diseño técnico escapa a las mismas fuerzas de economía política, presión sectorial y compromiso administrativo que han atravesado, sin excepción, cada una de las figuras que hemos construido juntos a lo largo de todo este tema.























- Cortes de Burgos de 1342 (origen de la alcabala castellana).
- Marqués de la Ensenada: Catastro de Ensenada (1749-1756) y proyecto de Única Contribución.
- Ley de Supresión del Impuesto de Consumos, de 12 de junio de 1911, y Real Decreto de 29 de junio de 1911.
- Ley de Reforma Tributaria de 1964 (reforma Navarro Rubio) y Texto Refundido del Impuesto General sobre el Tráfico de las Empresas.
- Ley 30/1985, de 2 de agosto, del Impuesto sobre el Valor Añadido (BOE de 9 de agosto de 1985), exposición de motivos.
- FUENTES QUINTANA, E.: propuestas de reforma de la imposición sobre ventas anteriores a 1964.
- FLORES DE LEMUS, A.: proyecto de reforma de la hacienda local durante la dictadura de Primo de Rivera (1923-1930).
- Ley 37/1992, de 28 de diciembre, del Impuesto sobre el Valor Añadido (ya citada).
- Ley 28/2014, de 27 de noviembre, por la que se modifica la Ley 37/1992 del IVA (introducción de la letra g) del art. 84.Uno.2º LIVA, con efectos desde el 1 de abril de 2015).
- Real Decreto 1073/2014, de 19 de diciembre, por el que se modifica el Reglamento del IVA (desarrollo de los nuevos supuestos de inversión del sujeto pasivo).
- Consulta Vinculante de la DGT V3274-15, de 26 de octubre de 2015 (concepto de empresario o profesional revendedor a efectos del art. 84.Uno.2º.g) LIVA).
- Real Decreto 596/2016, de 2 de diciembre, para la modernización, mejora e impulso del uso de medios electrónicos en la gestión del IVA (creación del SII).
- Ley 22/2009, de 18 de diciembre, por la que se regula el sistema de financiación de las CCAA de régimen común (ya citada en el documento original; ahora en su contenido específico sobre el índice de consumo territorial).
- Ley del Concierto Económico con el País Vasco y Ley del Convenio Económico con Navarra (límites de la capacidad normativa foral en materia de IVA por razón de la armonización comunitaria).
- Reglamento (UE) n.º 904/2010 del Consejo, de 7 de octubre de 2010, relativo a la cooperación administrativa y la lucha contra el fraude en el ámbito del IVA (creación de Eurofisc).
- Europol y Tribunal de Cuentas Europeo: estimaciones sobre concentración del fraude intracomunitario de IVA en manos de organizaciones criminales.
- Propuesta de Reglamento del Consejo sobre acceso de la Fiscalía Europea y la OLAF a la información del Reglamento (UE) n.º 904/2010 (desarrollo normativo más reciente, 2025-2026).
- Ley 22/2009, de 18 de diciembre, por la que se regula el sistema de financiación de las CCAA de régimen común (ya citada en el documento original; ahora en su contenido específico sobre el índice de consumo territorial).
- Ley del Concierto Económico con el País Vasco y Ley del Convenio Económico con Navarra (límites de la capacidad normativa foral en materia de IVA por razón de la armonización comunitaria).
- Reglamento (UE) n.º 904/2010 del Consejo, de 7 de octubre de 2010, relativo a la cooperación administrativa y la lucha contra el fraude en el ámbito del IVA (creación de Eurofisc).
- Europol y Tribunal de Cuentas Europeo: estimaciones sobre concentración del fraude intracomunitario de IVA en manos de organizaciones criminales.
- Propuesta de Reglamento del Consejo sobre acceso de la Fiscalía Europea y la OLAF a la información del Reglamento (UE) n.º 904/2010 (desarrollo normativo más reciente, 2025-2026).
- Ley 37/1992, de 28 de diciembre, del IVA, arts. 68 a 76 (lugar de realización del hecho imponible), arts. 75-77 (devengo), art. 78 (base imponible), art. 96 (exclusiones y restricciones del derecho a deducir).
- Real Decreto 1624/1992, de 29 de diciembre, Reglamento del IVA (desarrollo reglamentario de la devolución de saldos a favor del sujeto pasivo).
- Sentencia del TJCE de 15 de junio de 1989, Stichting Uitvoering Financiële Acties, C-348/87 (doctrina fundacional de interpretación estricta de las exenciones de IVA).
- Sentencia del TJUE de 7 de marzo de 2017, Comisión/Polonia, C-390/15, Gran Sala (tipo reducido y libros electrónicos).
- Directiva (UE) 2018/1713 del Consejo, de 6 de noviembre de 2018 (reforma del Anexo III de la Directiva 2006/112/CE para incluir publicaciones electrónicas).
- Sentencias del TJUE de 5 de marzo de 2020, C-513/20 y C-48/19 (interpretación estricta de las exenciones sanitarias).
- Sentencia del TJUE de 18 de septiembre de 2019, C-42/18 (interpretación estricta de las exenciones financieras).
- Real Decreto 1624/1992, de 29 de diciembre, exposición de motivos (referencia a las exenciones como "beneficios fiscales" adaptados a la normativa comunitaria).
- Anexo III de la Directiva 2006/112/CE (lista cerrada de categorías susceptibles de tipo reducido).
- Ley 37/1992, de 28 de diciembre, del IVA, art. 20 (exenciones en operaciones interiores) y art. 20.Dos (renuncia a la exención inmobiliaria).
- Real Decreto 1624/1992, de 29 de diciembre, Reglamento del IVA, art. 8 (requisitos formales de la renuncia a las exenciones inmobiliarias).
- Ley 37/1992, arts. 102-106 (regla de prorrata general y especial).
- Fondo Monetario Internacional (2010): A Fair and Substantial Contribution by the Financial Sector (ya citado en el Tema 13, aplicado ahora a la exención financiera española).
- Ley 37/1992, de 28 de diciembre, del IVA, arts. 90-91 (tipos de gravamen) y Anexo (relación de bienes y servicios a tipo reducido y superreducido).
- Real Decreto-ley 20/2022, de 27 de diciembre, de medidas de respuesta a las consecuencias económicas y sociales de la guerra en Ucrania (introducción del tipo del 0%).
- Real Decreto-ley 4/2024, de 26 de junio, por el que se prorrogan determinadas medidas para afrontar las consecuencias de los conflictos en Ucrania y Oriente Próximo (calendario de retirada gradual del tipo del 0%, e incorporación permanente del aceite de oliva al tipo superreducido).
- Ley 38/1992, de 28 de diciembre, de Impuestos Especiales.
- Real Decreto 1165/1995, de 7 de julio, Reglamento de los Impuestos Especiales.
- GRUBER, J. y KÖSZEGI, B. (2001): "Is Addiction 'Rational'? Theory and Evidence", Quarterly Journal of Economics, vol. 116, núm. 4.
- BERNHEIM, B. D. y RANGEL, A. (2004): "Addiction and Cue-Triggered Decision Processes", American Economic Review, vol. 94, núm. 5.
- O'DONOGHUE, T. y RABIN, M. (2006): "Optimal Sin Taxes", Journal of Public Economics, vol. 90, núms. 10-11.
- Ley 38/1992, de 28 de diciembre, de Impuestos Especiales
- Directiva 92/83/CEE del Consejo, de 19 de octubre de 1992, relativa a la armonización de las estructuras de los impuestos especiales sobre el alcohol y las bebidas alcohólicas.
- Directiva 92/84/CEE del Consejo, de 19 de octubre de 1992, relativa a la aproximación de los tipos del impuesto especial sobre el alcohol y las bebidas alcohólicas (tipo mínimo cero para el vino).
- Organización Mundial de la Salud: Global Strategy to Reduce the Harmful Use of Alcohol
- Alcohol (Minimum Pricing) (Scotland) Act 2012, en vigor desde el 1 de mayo de 2018.
- Ley 24/2001, de 27 de diciembre, de Medidas Fiscales, Administrativas y del Orden Social, art. 9 (creación del IVMDH).
- Ley 53/2002, de 30 de diciembre, de Medidas Fiscales, Administrativas y del Orden Social, art. 7 (modificación del IVMDH).
- Sentencia del TJUE de 27 de febrero de 2014, *Transportes Jordi Besora, S.L. contra Generalitat de Catalunya*, asunto C-82/12.
- Directiva 92/12/CEE del Consejo, de 25 de febrero de 1992, relativa al régimen general, tenencia, circulación y controles de los productos objeto de impuestos especiales.
- Analistas Financieros Internacionales (AFI) (2013): informe sobre comercio ilícito de tabaco en España, mayo de 2013.
- KPMG: *Project SUN* (informes anuales sobre comercio ilícito de tabaco en la Unión Europea, Suiza y Noruega, encargados por Philip Morris International).
- INGELMO, R.: tesis doctoral *"La incidencia de los IIEE sobre las labores del tabaco en la evolución del comercio ilícito de estos productos. Caso español"*.
- STOP (Stopping Tobacco Organizations and Products) / exposetobacco.org: análisis crítico de la metodología de los informes Project SUN y Project Stella.
- Directiva 2003/96/CE del Consejo, de 27 de octubre de 2003, por la que se reestructura el régimen comunitario de imposición de los productos energéticos y de la electricidad.
- Ley 15/2012, de 27 de diciembre, de medidas fiscales para la sostenibilidad energética (creación del IVPEE y de los impuestos sobre combustible nuclear gastado y residuos radiactivos).
- Sentencia del TJUE de 3 de marzo de 2021, *Promociones Oliva Park, S.L.*, asunto C-220/19 (compatibilidad del IVPEE con el Derecho de la Unión).
- ORTIZ CALLE, E. (2019): "Compatibilidad del impuesto sobre el valor de la producción de la energía eléctrica con el derecho de la Unión Europea", *Revista de Contabilidad y Tributación. CEF*, núms. 437-438.
- Ley 38/1992, de 28 de diciembre, de Impuestos Especiales, Título II (impuesto especial sobre determinados medios de transporte), exposición de motivos.
- Ley 34/2007, de 15 de noviembre, de calidad del aire y protección de la atmósfera, disposición final octava (reestructuración del IEDMT según emisiones de CO2).
- Propuesta de Directiva del Consejo relativa a los impuestos sobre turismos, COM(2005) 261 final, retirada por la Comisión Europea en 2012.
- Sentencia del TJUE de 2 de septiembre de 2021, *Comisión/Portugal (Impuesto que grava los vehículos)*, asunto C-169/20.
- Sentencia del TJUE de 22 de febrero de 2024, *Comisión/Malta*, asunto C-694/22.
- Agencia Estatal de Administración Tributaria: *Estadística de Matriculación de Vehículos*, ejercicio 2008.







































































\clearpage
\section{Conclusión} 


\subsection{Recapitulación}


\subsection{Extensiones y relación con otras partes del Temario}



\subsection{Opinión}

\subsubsection{Idea final (Salida o cierra)}

\clearpage
\pagenumbering{gobble}
\MyBibliography



\MyBibItem{DAWID y GATTI}{Chapter 2 - Agent-Based Macroeconomics}{2018}{https://doi.org/10.1016/bs.hescom.2018.02.006}


% ANEXOS
\clearpage










\end{document}